\documentclass[11pt]{article}

\usepackage[utf8]{inputenc}
\usepackage[T1]{fontenc}
\usepackage[ngerman]{babel}
\usepackage{amsmath, amssymb, amsthm}
\usepackage{bm}
\usepackage{geometry}
\usepackage{booktabs}
\usepackage{array}
\usepackage{xcolor}
\usepackage{hyperref}
\usepackage{microtype}
\usepackage{parskip}
\usepackage{enumitem}
\usepackage{tcolorbox}
\tcbuselibrary{skins, breakable}
\usepackage{mdframed}
\usepackage{pgfplots}
\pgfplotsset{compat=1.18}
\usepackage{tikz}
\usepackage{mathtools}
\usepackage{bbm}

\geometry{left=2.8cm, right=2.8cm, top=3cm, bottom=3cm}

% ------------------------------------------------------------------
% Farben (angelehnt an das Skript)
% ------------------------------------------------------------------
\definecolor{mainblue}{RGB}{0, 83, 156}
\definecolor{lightblue}{RGB}{230, 242, 255}
\definecolor{algbg}{RGB}{245, 248, 255}
\definecolor{defbg}{RGB}{255, 252, 235}
\definecolor{thmbg}{RGB}{240, 255, 240}
\definecolor{webbg}{RGB}{255, 244, 230}
\definecolor{srcbg}{RGB}{240, 240, 240}

% ------------------------------------------------------------------
% Theorem-Umgebungen
% ------------------------------------------------------------------
\newtheoremstyle{csthm}{}{}{\itshape}{}{\bfseries\color{mainblue}}{.}{ }{}
\theoremstyle{csthm}
\newtheorem{theorem}{Satz}[section]
\newtheorem{lemma}[theorem]{Lemma}
\newtheorem{corollary}[theorem]{Korollar}

\newtheoremstyle{csdef}{}{}{}{}{\bfseries\color{mainblue}}{.}{ }{}
\theoremstyle{csdef}
\newtheorem{definition}[theorem]{Definition}
\newtheorem{remark}[theorem]{Bemerkung}

% ------------------------------------------------------------------
% Algorithmus-Box
% ------------------------------------------------------------------
\tcbset{
  algostyle/.style={
    enhanced, colback=algbg, colframe=mainblue!60,
    fonttitle=\bfseries\color{mainblue},
    title={Algorithmus}, breakable, top=3mm, bottom=3mm,
  },
  defstyle/.style={
    enhanced, colback=defbg, colframe=mainblue!40,
    breakable, top=2mm, bottom=2mm, left=4mm,
  },
  thmstyle/.style={
    enhanced, colback=thmbg, colframe=mainblue!50,
    breakable, top=2mm, bottom=2mm, left=4mm,
  },
  webstyle/.style={
    enhanced, colback=webbg, colframe=mainblue!30!orange,
    fonttitle=\bfseries\color{mainblue!30!orange!80!black},
    title={Zusatzinfo (Internet)}, breakable, top=2mm, bottom=2mm, left=4mm,
  },
  srcstyle/.style={
    enhanced, colback=srcbg, colframe=gray!60,
    fonttitle=\bfseries\itshape\color{gray!60!black},
    title={Quelle}, breakable, top=1mm, bottom=1mm, left=3mm, boxrule=0.4pt,
    fontupper=\small\itshape,
  }
}

% ------------------------------------------------------------------
% Makros
% ------------------------------------------------------------------
\newcommand{\R}{\mathbb{R}}
\newcommand{\E}{\mathbb{E}}
\newcommand{\I}{\mathbbm{1}}
\newcommand{\N}{\mathcal{N}}
\newcommand{\X}{\mathcal{X}}
\newcommand{\Y}{\mathcal{Y}}
\newcommand{\thetahat}{\hat{\theta}}
\newcommand{\htheta}{h_\theta}
\newcommand{\norm}[1]{\left\|#1\right\|}
\newcommand{\abs}[1]{\left|#1\right|}
\newcommand{\T}{^\top}
\newcommand{\indep}{\perp\!\!\!\perp}

\hypersetup{colorlinks=true, linkcolor=mainblue, urlcolor=mainblue, citecolor=mainblue}

% ------------------------------------------------------------------
% Titelseite
% ------------------------------------------------------------------
\title{%
  \vspace{-1cm}
  {\Large CS229: Machine Learning}\\[6pt]
  {\LARGE\bfseries\color{mainblue} Vorlesungsskript -- Zusammenfassung}\\[6pt]
  {\large Andrew Ng \& Tengyu Ma}
}
\author{Stanford University}
\date{Ausgabe 2023 -- Deutsche Zusammenfassung}

% ==================================================================
\begin{document}
% ==================================================================
% ngerman/babel macht " zu einem aktiven Kürzel-Zeichen (z.B. "a -> ä,
% "s -> ß). Da " im Text auch als normales Anführungszeichen genutzt wird,
% deaktivieren wir dieses Shorthand global, um Fehlsubstitutionen zu vermeiden.
\shorthandoff{"}
\maketitle
\tableofcontents
\newpage

% ==================================================================
\section*{Vorwort}
% ==================================================================

Diese Zusammenfassung begleitet die CS229-Vorlesungsreihe -- basierend auf
dem Skript \texttt{main\_notes.pdf}, sämtlichen ergänzenden Lecture Notes im
Ordner \texttt{notes/} und \texttt{extra-notes/} (Decision Trees, Ensembling,
Perceptron, Faktorenanalyse, Fehleranalyse sowie John Duchis Supplemental
Notes zu Hoeffdings Ungleichung, Verlustfunktionen und dem
Representer-Theorem) sowie zusätzlichem Material aus \texttt{materials/} und
\texttt{section/} (Andrew Ngs ML-Advice-Folien, Konvexe Optimierung/KKT,
SMO-Pseudocode, Gaussian Processes, Hidden Markov Models, Evaluationsmetriken
sowie kritische Perspektiven auf ML). Ziel ist es, die wesentlichen Konzepte,
Algorithmen und mathematischen Grundlagen aller 27 Kapitel kompakt
darzustellen -- mit Herleitungen, Intuition und Querverweisen zwischen den
Themen. Zwei Arten von Boxen kennzeichnen Material außerhalb des
main\_notes-Kerntexts: graue \textbf{Quelle}-Boxen verweisen auf zusätzliche,
aber offizielle CS229-Materialien (Section Notes, Handouts) aus diesem
Repository; orange \textbf{Zusatzinfo (Internet)}-Boxen enthalten recherchiertes
Material, das nicht aus CS229-Unterlagen stammt.

Das Skript ist in fünf Teile gegliedert:
\begin{enumerate}
  \item \textbf{Überwachtes Lernen} (Kap.~1--10): Regression, Klassifikation,
        generative Modelle, Kernel-Methoden, SVMs, Perceptron, Decision
        Trees, Ensemble-Methoden, Verlustfunktionen
  \item \textbf{Deep Learning} (Kap.~11): Neuronale Netze und Backpropagation
  \item \textbf{Generalisierung \& Regularisierung} (Kap.~12--16):
        Bias-Varianz, Double Descent, Modellauswahl, Gaussian Processes,
        praktische Fehleranalyse, Evaluationsmetriken
  \item \textbf{Unüberwachtes Lernen} (Kap.~17--24): Clustering, EM,
        Faktorenanalyse, Hidden Markov Models, PCA, ICA, Foundation Models,
        kritische Perspektiven auf ML
  \item \textbf{Bestärkendes Lernen} (Kap.~25--27): MDPs, LQR, Policy
        Gradient
\end{enumerate}

\noindent\textbf{Notation.} Vektoren und Matrizen werden fett geschrieben
($\bm{x}, \bm{W}$). Skalare kursiv ($x, \theta$). Zufallsgrößen mit
Großbuchstaben ($X, Y$). Trainingsdaten $\{(x^{(i)}, y^{(i)})\}_{i=1}^n$ mit
$n$ Beispielen.

\newpage

% ==================================================================
\part{Überwachtes Lernen}
% ==================================================================

% ------------------------------------------------------------------
\section{Lineare Regression}
% ------------------------------------------------------------------

\subsection{Aufgabenstellung}

Beim überwachten Lernen sei ein Datensatz
$\{(x^{(i)}, y^{(i)})\}_{i=1}^n$ mit $x^{(i)} \in \R^d$ und
$y^{(i)} \in \R$ gegeben. Ziel ist es, eine Hypothesenfunktion
$h : \R^d \to \R$ zu lernen, die für neue Eingaben $x$ gute
Vorhersagen $h(x) \approx y$ liefert.

Bei der \textbf{linearen Regression} wird angesetzt:
\begin{equation}
  \htheta(x) = \theta\T x = \sum_{j=0}^{d} \theta_j x_j,
  \qquad x_0 := 1 \text{ (Bias-Term).}
\end{equation}

\subsection{Kostenfunktion und Gradientenverfahren}

Die \textbf{mittlere quadratische Abweichung} (MSE) dient als Verlustfunktion:
\begin{equation}
  J(\theta) = \frac{1}{2n} \sum_{i=1}^{n}
    \bigl(\htheta(x^{(i)}) - y^{(i)}\bigr)^2.
\end{equation}

Batch Gradient Descent minimiert $J$ durch iterative Updates:
\begin{equation}
  \theta_j \leftarrow \theta_j
    - \alpha \frac{\partial J}{\partial \theta_j}
  = \theta_j + \frac{\alpha}{n}
    \sum_{i=1}^{n} \bigl(y^{(i)} - \htheta(x^{(i)})\bigr)\, x_j^{(i)}.
\end{equation}

\noindent\textbf{Stochastisches Gradientenverfahren (SGD).} Bei großen
Datensätzen ist ein vollständiger Batch-Update teuer. SGD wählt in jedem
Schritt ein einzelnes (oder eine kleine Gruppe von) Beispiel(en) und
approximiert den Gradienten:
\begin{equation}
  \theta_j \leftarrow \theta_j
    + \alpha \bigl(y^{(i)} - \htheta(x^{(i)})\bigr)\, x_j^{(i)}.
\end{equation}

\subsection{Normalgleichungen (Geschlossene Lösung)}

Die \textbf{Design-Matrix} $X \in \R^{n \times d}$ enthält die
Trainingsvektoren zeilenweise. Dann lässt sich der Verlust kompakt
schreiben als
\begin{equation}
  J(\theta) = \tfrac{1}{2n}\, \norm{X\theta - \bm{y}}^2.
\end{equation}
Ableiten und Null setzen ergibt die \textbf{Normalgleichungen}:
\begin{equation}
  \boxed{\thetahat = (X\T X)^{-1} X\T \bm{y}.}
\end{equation}
Diese geschlossene Lösung ist für kleine bis mittlere $d$ effizient; für
sehr große $d$ oder schlecht konditioniertes $X\T X$ wird das iterative
Verfahren bevorzugt.

\subsection{Probabilistische Interpretation}

Annahme: $y^{(i)} = \theta\T x^{(i)} + \varepsilon^{(i)}$ mit iid
$\varepsilon^{(i)} \sim \N(0, \sigma^2)$. Dann ist die Likelihood
\[
  L(\theta) = \prod_{i=1}^n p(y^{(i)} \mid x^{(i)}; \theta)
  = \prod_{i=1}^n \frac{1}{\sqrt{2\pi}\,\sigma}
    \exp\!\Bigl(-\frac{(y^{(i)} - \theta\T x^{(i)})^2}{2\sigma^2}\Bigr).
\]
Maximierung der Log-Likelihood ist äquivalent zur Minimierung von $J(\theta)$
-- die Least-Squares-Methode ist somit der Maximum-Likelihood-Schätzer (MLE)
unter der Gauß-Annahme.

\subsection{Lokal gewichtete lineare Regression}

Bei der \textbf{lokal gewichteten Regression} (LWLR) werden die
Trainingspunkte je nach Nähe zur Anfrage $x$ unterschiedlich gewichtet:
\begin{equation}
  w^{(i)}(x) = \exp\!\Bigl(-\frac{\norm{x^{(i)} - x}^2}{2\tau^2}\Bigr),
\end{equation}
und der gewichtete Verlust $J(\theta) = \sum_i w^{(i)}(x)(y^{(i)} -
\theta\T x^{(i)})^2$ wird minimiert. LWLR ist ein
\textbf{nichtparametrisches} Verfahren: die Lösung hängt stets von den
Trainingsdaten ab.

% ------------------------------------------------------------------
\section{Klassifikation und Logistische Regression}
% ------------------------------------------------------------------

\subsection{Logistische Regression}

Bei binärer Klassifikation ($y \in \{0,1\}$) ist ein lineares Modell
ungeeignet, da $\theta\T x \notin [0,1]$. Stattdessen setzt man die
\textbf{Sigmoid-Funktion} ein:
\begin{equation}
  g(z) = \frac{1}{1 + e^{-z}}, \qquad
  \htheta(x) = g(\theta\T x) = \frac{1}{1 + e^{-\theta\T x}}.
\end{equation}

Das stochastische Modell lautet:
\begin{align}
  P(y = 1 \mid x;\,\theta) &= \htheta(x), \\
  P(y = 0 \mid x;\,\theta) &= 1 - \htheta(x).
\end{align}

Die Log-Likelihood über den Datensatz ist
\begin{equation}
  \ell(\theta) = \sum_{i=1}^n
    \Bigl[y^{(i)}\log \htheta(x^{(i)})
      + (1 - y^{(i)})\log\bigl(1 - \htheta(x^{(i)})\bigr)\Bigr].
\end{equation}
Maximierung via Stochastic Gradient Ascent führt auf die Update-Regel
\begin{equation}
  \theta_j \leftarrow \theta_j
    + \alpha\bigl(y^{(i)} - \htheta(x^{(i)})\bigr)\, x_j^{(i)},
\end{equation}
die formal identisch mit der LMS-Regel aussieht -- ein bemerkenswertes
Ergebnis, obwohl $h$ eine völlig andere Funktion ist.

\subsection{Newton-Verfahren (Fisher Scoring)}

Das Newton-Raphson-Verfahren konvergiert deutlich schneller als
Gradientenverfahren, weil es Krümmungsinformation nutzt:
\begin{equation}
  \theta \leftarrow \theta - H^{-1} \nabla_\theta \ell(\theta),
\end{equation}
wobei die \textbf{Hesse-Matrix} $H_{jk} = \partial^2\ell / (\partial\theta_j
\partial\theta_k)$ für logistische Regression lautet:
\begin{equation}
  H = -\frac{1}{n} X\T S\, X, \qquad
  S = \operatorname{diag}\bigl(\htheta(x^{(1)})(1-\htheta(x^{(1)})),\ldots\bigr).
\end{equation}
Newton-Schritt heißt hier auch \textbf{Fisher Scoring} und konvergiert
quadratisch -- ein Schritt entspricht grob $O(d)$ Gradientenschritten.

\subsection{Mehrklassen-Klassifikation und Softmax}

Für $K$ Klassen parametrisiert man für jede Klasse $k$ einen Vektor
$\theta_k \in \R^d$ und nutzt die \textbf{Softmax}-Funktion:
\begin{equation}
  P(y = k \mid x;\,\Theta)
  = \frac{\exp(\theta_k\T x)}{\sum_{j=1}^K \exp(\theta_j\T x)}.
\end{equation}
Die zugehörige Verlustfunktion ist der \textbf{Cross-Entropy-Verlust}:
\begin{equation}
  \ell_\mathrm{CE}(\hat{y}, y)
  = -\log P(y \mid x;\,\Theta)
  = -\theta_y\T x + \log\sum_{j=1}^K \exp(\theta_j\T x).
\end{equation}

% ------------------------------------------------------------------
\section{Generalisierte Lineare Modelle (GLMs)}
% ------------------------------------------------------------------

\subsection{Die Exponentialfamilie}

Viele Standardverteilungen lassen sich in der \textbf{Exponentialfamilien}-
Form schreiben:
\begin{equation}
  p(y;\,\eta) = b(y)\exp\!\bigl(\eta\T T(y) - a(\eta)\bigr),
\end{equation}
mit natürlichem Parameter $\eta$, hinreichender Statistik $T(y)$,
Log-Partitionsfunktion $a(\eta)$ und Basismaß $b(y)$.

\begin{tcolorbox}[defstyle]
\textbf{Bernoulli-Verteilung:} $p(y;\phi) = \phi^y(1-\phi)^{1-y}$,
$y\in\{0,1\}$. Schreibweise als Exponentialfamilie mit
$\eta = \log\frac{\phi}{1-\phi}$, also $\phi = 1/(1+e^{-\eta})$ -- dies
\emph{ist} genau die Sigmoid-Funktion.

\medskip
\textbf{Gauß-Verteilung:} $\N(\mu, \sigma^2)$ mit $\eta = \mu$, $T(y) = y$,
$a(\eta) = \eta^2/2$, $b(y) = (2\pi)^{-1/2}e^{-y^2/2}$.
\end{tcolorbox}

\subsection{GLM-Konstruktion}

Ein \textbf{Generalisiertes Lineares Modell} (GLM) basiert auf drei Annahmen:
\begin{enumerate}
  \item $y \mid x;\,\theta \sim \text{ExponentialFamily}(\eta)$.
  \item Hypothese: $h(x) = \E[T(y) \mid x;\,\theta]$.
  \item Lineare Verbindung: $\eta = \theta\T x$.
\end{enumerate}

Die kanonische \textbf{Responsefunktion} $g(\eta) = \E[T(y);\,\eta]$
verbindet natürlichen Parameter und Erwartungswert. Ihre Inverse ist die
\textbf{Linkfunktion}.

\begin{center}
\begin{tabular}{llll}
\toprule
Verteilung & Responsefunktion & Modell & Anwendung \\
\midrule
Gauß & Identität & Lineare Regression & Regressionsaufgaben \\
Bernoulli & Sigmoid & Logistische Regression & Binäre Klassifikation \\
Multinomial & Softmax & Softmax-Regression & Mehrklassenprobleme \\
Poisson & $e^\eta$ & Poisson-Regression & Zähldaten \\
\bottomrule
\end{tabular}
\end{center}

% ------------------------------------------------------------------
\section{Generative Lernalgorithmen}
% ------------------------------------------------------------------

\subsection{Diskriminativ vs. Generativ}

\textbf{Diskriminative} Modelle (z.\,B. logistische Regression) modellieren
$p(y \mid x)$ direkt. \textbf{Generative} Modelle lernen $p(x \mid y)$ und
$p(y)$ und nutzen den Satz von Bayes:
\begin{equation}
  p(y \mid x) = \frac{p(x \mid y)\,p(y)}{p(x)},
  \qquad p(x) = \sum_y p(x \mid y)\,p(y).
\end{equation}

\subsection{Gaussian Discriminant Analysis (GDA)}

GDA modelliert $x \mid y$ als multivariate Normalverteilung:
\begin{align}
  y &\sim \operatorname{Bernoulli}(\phi), \\
  x \mid y=0 &\sim \N(\mu_0, \Sigma), \\
  x \mid y=1 &\sim \N(\mu_1, \Sigma).
\end{align}
Die \textbf{multivariate Normalverteilung} lautet:
\begin{equation}
  p(x;\,\mu,\Sigma)
  = \frac{1}{(2\pi)^{d/2}|\Sigma|^{1/2}}
    \exp\!\Bigl(-\tfrac{1}{2}(x-\mu)\T\Sigma^{-1}(x-\mu)\Bigr).
\end{equation}

Die MLE-Schätzer sind:
\begin{align}
  \hat\phi &= \frac{1}{n}\sum_{i=1}^n \I\{y^{(i)}=1\}, \\
  \hat\mu_k &= \frac{\sum_i \I\{y^{(i)}=k\}\,x^{(i)}}
                    {\sum_i \I\{y^{(i)}=k\}},
  \quad k \in \{0,1\}, \\
  \hat\Sigma &= \frac{1}{n}\sum_{i=1}^n
    (x^{(i)} - \hat\mu_{y^{(i)}})(x^{(i)} - \hat\mu_{y^{(i)}})\T.
\end{align}

\begin{tcolorbox}[thmstyle]
\textbf{GDA und logistische Regression.}
Wenn $p(x \mid y)$ tatsächlich Gauß-verteilt ist, hat die Posterior
$p(y=1 \mid x)$ stets die logistische Form $1/(1+e^{-\theta\T x})$.
GDA macht \emph{stärkere} Annahmen und ist asymptotisch effizienter, wenn
diese zutreffen. Logistische Regression ist robuster, wenn die
Gauß-Annahme verletzt ist.
\end{tcolorbox}

\subsection{Naive Bayes}

Naive Bayes setzt \textbf{bedingte Unabhängigkeit} der Features voraus:
\begin{equation}
  p(x \mid y) = \prod_{j=1}^d p(x_j \mid y).
\end{equation}
Dies vereinfacht die Schätzung erheblich. Für Textklassifikation mit
Binärfeatures:
\begin{equation}
  \hat\phi_{j\mid y=k}
  = \frac{\sum_i \I\{x_j^{(i)}=1,\, y^{(i)}=k\}}
         {\sum_i \I\{y^{(i)}=k\}}.
\end{equation}

\textbf{Laplace-Glättung} verhindert Nullwahrscheinlichkeiten für
ungesehene Features:
\begin{equation}
  \hat\phi_{j\mid y=k}
  = \frac{1 + \sum_i \I\{x_j^{(i)}=1,\, y^{(i)}=k\}}
         {2 + \sum_i \I\{y^{(i)}=k\}}.
\end{equation}

% ------------------------------------------------------------------
\section{Kernel-Methoden}
% ------------------------------------------------------------------

\subsection{Feature Maps und das Kernel-Trick}

Sei $\phi : \R^d \to \R^p$ eine Feature-Abbildung (z.\,B. Polynom-Features).
Das LMS-Update in Feature-Raum lautet:
\[
  \theta \leftarrow \theta
    + \alpha\sum_{i=1}^n \bigl(y^{(i)} - \theta\T\phi(x^{(i)})\bigr)\phi(x^{(i)}).
\]
Da $\theta$ stets als Linearkombination der $\phi(x^{(i)})$ geschrieben werden
kann ($\theta = \sum_i \beta_i \phi(x^{(i)})$), reduziert sich die Vorhersage
auf Skalarprodukte:
\begin{equation}
  h(x) = \theta\T\phi(x) = \sum_{i=1}^n \beta_i\,\phi(x^{(i)})\T\phi(x)
        = \sum_{i=1}^n \beta_i\, K(x^{(i)}, x).
\end{equation}

\begin{definition}
  Eine Funktion $K : \R^d \times \R^d \to \R$ heißt \textbf{Kernel}, wenn
  $K(x, z) = \phi(x)\T\phi(z)$ für eine Feature-Abbildung $\phi$.
\end{definition}

Der entscheidende Vorteil: $K(x,z)$ kann in $O(d)$ Zeit berechnet werden,
obwohl $\phi(x)$ möglicherweise sehr hochdimensional (oder unendlich-dimensional)
ist. Beispiel: Polynom-Kernel $K(x,z) = (x\T z + c)^k$ entspricht einem
$O(d^k)$-dimensionalen Feature-Raum.

\subsection{Wichtige Kernels}

\begin{center}
\begin{tabular}{lll}
\toprule
Kernel & $K(x, z)$ & Feature-Raum \\
\midrule
Linear & $x\T z$ & Original-Raum $\R^d$ \\
Polynom (Grad $k$) & $(x\T z + c)^k$ & $O(d^k)$-dimensional \\
Gauß (RBF) & $\exp\!\bigl(-\frac{\norm{x-z}^2}{2\sigma^2}\bigr)$ & $\infty$-dimensional \\
\bottomrule
\end{tabular}
\end{center}

\subsection{Mercers Satz}

\begin{tcolorbox}[thmstyle]
\textbf{Mercers Satz.}
Eine Funktion $K : \R^d \times \R^d \to \R$ ist genau dann ein gültiger
Kernel (d.\,h. $K(x,z) = \phi(x)\T\phi(z)$ für ein $\phi$), wenn für beliebige
Punkte $x^{(1)},\ldots,x^{(n)}$ die \textbf{Gram-Matrix}
$[K]_{ij} = K(x^{(i)}, x^{(j)})$ symmetrisch und positiv semidefinit ist.
\end{tcolorbox}

Kernels sind unter Addition, positiver Skalierung und (elementweisem) Produkt
abgeschlossen -- diese Rechenregeln ermöglichen den systematischen Aufbau
komplexer Kernels aus einfachen.

\subsection{Representer-Theorem}

Warum genügt es, $\theta$ stets als Linearkombination der Trainingsbeispiele
anzusetzen? Für einen regularisierten empirischen Risiko-Ansatz
\begin{equation}
  J_\lambda(\theta) = \frac{1}{n}\sum_{i=1}^n L(\theta\T x^{(i)}, y^{(i)})
    + \frac{\lambda}{2}\norm{\theta}^2,
\end{equation}
mit einem in $z$ konvexen Verlust $L(z,y)$ (z.\,B.\ quadratischer, logistischer
oder Hinge-Verlust) und $\lambda \geq 0$ gilt:

\begin{tcolorbox}[thmstyle]
\textbf{Representer-Theorem.} Es existiert ein Minimierer von $J_\lambda$,
der sich als Linearkombination der Trainingsbeispiele schreiben lässt:
\begin{equation}
  \theta = \sum_{i=1}^n \alpha_i\, x^{(i)}
\end{equation}
für reelle Koeffizienten $\alpha_1,\ldots,\alpha_n$.
\end{tcolorbox}

\textbf{Beweisidee} (differenzierbarer Fall, $\lambda>0$): Am Minimum
verschwindet der Gradient,
\[
  \nabla J_\lambda(\theta) = \frac1n\sum_{i=1}^n L'(\theta\T x^{(i)},y^{(i)})\,x^{(i)}
    + \lambda\theta = 0.
\]
Auflösen nach $\theta$ liefert direkt
$\theta = -\frac{1}{n\lambda}\sum_i L'(\theta\T x^{(i)},y^{(i)})\,x^{(i)}
=: \sum_i \alpha_i x^{(i)}$.
Dies rechtfertigt im Nachhinein den Kernel-Trick aus Abschnitt~5.1: Da die
Lösung stets in der von den Daten aufgespannten Linearkombination liegt,
genügen zur Auswertung der Hypothese stets nur Skalarprodukte
$\phi(x^{(i)})\T\phi(x^{(j)}) = K(x^{(i)}, x^{(j)})$ -- eine explizite
Darstellung von $\theta$ im (ggf.\ unendlich-dimensionalen) Feature-Raum ist
nie nötig.

% ------------------------------------------------------------------
\section{Support Vector Machines (SVMs)}
% ------------------------------------------------------------------

\subsection{Margins und Intuition}

Eine Hyperebene $\{x : w\T x + b = 0\}$ klassifiziert Punkt $x^{(i)}$
korrekt, wenn $y^{(i)}(w\T x^{(i)} + b) > 0$. Der \textbf{funktionale
Margin} misst diese Konfidenz:
\begin{equation}
  \hat\gamma^{(i)} = y^{(i)}(w\T x^{(i)} + b),
\end{equation}
und der \textbf{geometrische Margin} (abstandsbasiert) ist
\begin{equation}
  \gamma^{(i)} = \frac{y^{(i)}(w\T x^{(i)} + b)}{\norm{w}}.
\end{equation}
Das SVM-Ziel: finde die Hyperebene, die alle Punkte korrekt klassifiziert
und deren \emph{minimaler geometrischer Margin} maximal ist.

\subsection{Primal-Problem (linear trennbar)}

Mit der Normierung $\min_i \hat\gamma^{(i)} = 1$ wird das
\textbf{Optimale-Margin-Klassifikatorproblem}:
\begin{equation}
  \min_{w,b}\; \tfrac{1}{2}\norm{w}^2 \quad
  \text{s.t.}\quad y^{(i)}(w\T x^{(i)} + b) \geq 1, \quad i=1,\ldots,n.
\end{equation}
Dies ist ein konvexes quadratisches Programm.

\subsection{Lagrange-Dualität und Dual-Form}

Der Lagrangian ist
\[
  \mathcal{L}(w,b,\alpha)
  = \tfrac{1}{2}\norm{w}^2 - \sum_{i=1}^n \alpha_i
    \bigl[y^{(i)}(w\T x^{(i)} + b) - 1\bigr].
\]
Aus den KKT-Stationaritätsbedingungen folgt $w = \sum_i \alpha_i y^{(i)} x^{(i)}$.
Einsetzen ergibt das \textbf{Dual-Problem}:
\begin{equation}
  \max_\alpha\; W(\alpha) = \sum_{i=1}^n \alpha_i
    - \tfrac{1}{2}\sum_{i,j} \alpha_i\alpha_j y^{(i)} y^{(j)} \langle x^{(i)}, x^{(j)}\rangle
\end{equation}
s.t.\ $\alpha_i \geq 0$ und $\sum_i \alpha_i y^{(i)} = 0$.

Die Vorhersage hängt damit nur noch von Skalarprodukten ab -- und lässt sich
mit einem Kernel verallgemeinern: $\langle x^{(i)}, x^{(j)}\rangle \to K(x^{(i)}, x^{(j)})$.

\textbf{Support Vectors} sind die Punkte mit $\alpha_i > 0$; nur sie
definieren die Entscheidungsgrenze.

Die "KKT-Stationaritätsbedingungen", die oben zur Herleitung von $w$ genutzt
wurden, sind ein Spezialfall eines allgemeinen Satzes der konvexen
Optimierung, der Optimalität für \emph{jedes} Primal-Dual-Paar
charakterisiert:

\begin{tcolorbox}[thmstyle]
\textbf{Karush-Kuhn-Tucker(KKT)-Theorem.} Für ein konvexes Problem
$\min_x f(x)$ s.\,t.\ $g_i(x)\leq 0\;(i=1,\ldots,m)$, $h_i(x)=0\;(i=1,\ldots,p)$
mit Lagrangian $\mathcal L(x,\alpha,\beta) = f(x) + \sum_i \alpha_i g_i(x) +
\sum_i \beta_i h_i(x)$ gilt: $x^*$ ist primal-optimal und $(\alpha^*,\beta^*)$
dual-optimal genau dann, wenn (bei starker Dualität) die folgenden vier
Bedingungen erfüllt sind:
\begin{enumerate}[leftmargin=*, itemsep=1pt]
  \item \textbf{Primale Zulässigkeit}: $g_i(x^*)\leq 0$, $h_i(x^*)=0$.
  \item \textbf{Duale Zulässigkeit}: $\alpha_i^* \geq 0$.
  \item \textbf{Komplementärer Schlupf (complementary slackness)}: $\alpha_i^*\,g_i(x^*)=0$.
  \item \textbf{Stationarität}: $\nabla_x \mathcal L(x^*,\alpha^*,\beta^*)=0$.
\end{enumerate}
\end{tcolorbox}

Bedingung~3 erklärt direkt, warum nur Support Vectors ($\alpha_i>0$) die
Lösung bestimmen: Für alle anderen Punkte muss die zugehörige
Ungleichungsnebenbedingung $g_i(x^*)=0$ sein (d.\,h.\ $y^{(i)}(w\T x^{(i)}+b)=1$,
der Punkt liegt exakt auf dem Margin) oder $\alpha_i^*=0$ (der Punkt trägt
nicht zur Lösung bei).

\begin{tcolorbox}[srcstyle]
Quelle: \emph{Convex Optimization Overview (cnt'd)}, Chuong B. Do (CS229
Section Notes) -- \texttt{section/cs229-cvxopt2.pdf}, Abschnitt~1.6.
\end{tcolorbox}

\subsection{Soft-Margin SVM (nicht-trennbarer Fall)}

Wenn die Daten nicht linear trennbar sind, erlaubt man Schlupfvariablen
$\xi_i \geq 0$ (Verletzungen der Margin-Bedingung):
\begin{equation}
  \min_{w,b,\xi}\; \tfrac{1}{2}\norm{w}^2 + C\sum_{i=1}^n \xi_i \quad
  \text{s.t.}\quad y^{(i)}(w\T x^{(i)} + b) \geq 1 - \xi_i,\quad \xi_i \geq 0.
\end{equation}
Der Hyperparameter $C > 0$ steuert den Trade-off zwischen großem Margin und
kleinen Fehlern. Das Dual-Problem ändert sich nur in $0 \leq \alpha_i \leq C$.

\subsection{SMO-Algorithmus}

Der \textbf{Sequential Minimal Optimization (SMO)} Algorithmus löst das
quadratische Programm, indem er iterativ immer zwei $\alpha_i$ simultan
aktualisiert (da die Gleichheitsbedingung $\sum_i \alpha_i y^{(i)} = 0$ ein
einzelnes Update verbietet). Die Konvergenz wird durch Prüfung der
KKT-Bedingungen überwacht. Für die Soft-Margin-SVM lauten die KKT-Bedingungen
konkret:
\begin{equation}
  \alpha_i = 0 \;\Rightarrow\; y^{(i)}(w\T x^{(i)}{+}b)\geq 1, \quad
  \alpha_i = C \;\Rightarrow\; y^{(i)}(w\T x^{(i)}{+}b)\leq 1, \quad
  0{<}\alpha_i{<}C \;\Rightarrow\; y^{(i)}(w\T x^{(i)}{+}b) = 1.
\end{equation}

\textbf{Vereinfachtes SMO} (wie in Aufgabenstellungen zu Problem Set~2
verwendet): Statt einer aufwendigen Heuristik zur Auswahl des Paares
$(\alpha_i,\alpha_j)$ iteriert man einfach über alle $\alpha_i$; verletzt
$\alpha_i$ die KKT-Bedingungen, wird ein zufälliges $\alpha_j$ gewählt und
beide gemeinsam optimiert:

\begin{tcolorbox}[algostyle, title={Update-Schritt für ein Paar $(\alpha_i,\alpha_j)$}]
\textbf{1. Grenzen bestimmen} (damit $0\leq\alpha_j\leq C$ erhalten bleibt):
\[
  y^{(i)}\neq y^{(j)}:\; L=\max(0,\alpha_j{-}\alpha_i),\; H=\min(C,C{+}\alpha_j{-}\alpha_i)
  \qquad
  y^{(i)}=y^{(j)}:\; L=\max(0,\alpha_i{+}\alpha_j{-}C),\; H=\min(C,\alpha_i{+}\alpha_j)
\]
\textbf{2. Unbeschränktes Optimum und Clipping} mit
$E_k = f(x^{(k)}) - y^{(k)}$ (Vorhersagefehler) und
$\eta = 2K(x^{(i)},x^{(j)}) - K(x^{(i)},x^{(i)}) - K(x^{(j)},x^{(j)})$:
\[
  \alpha_j \leftarrow \alpha_j - \frac{y^{(j)}(E_i-E_j)}{\eta}, \qquad
  \alpha_j \leftarrow \operatorname{clip}(\alpha_j,\,L,\,H).
\]
\textbf{3. Update von $\alpha_i$}:
$\alpha_i \leftarrow \alpha_i + y^{(i)}y^{(j)}(\alpha_j^{\text{alt}}-\alpha_j)$.

\textbf{4. Bias $b$ neu setzen} so, dass die KKT-Bedingungen für $i,j$ erfüllt
bleiben (Mittelwert der beiden Randfälle, falls beide $\alpha$ im Inneren
$(0,C)$ liegen).
\end{tcolorbox}

\begin{tcolorbox}[srcstyle]
Quelle: \emph{CS229 Simplified SMO Algorithm}, CS229 Autumn 2009 (basierend
auf J. Platts Originalarbeit) -- \texttt{materials/smo.pdf}.
\end{tcolorbox}

% ------------------------------------------------------------------
\section{Perceptron und Online Learning}
% ------------------------------------------------------------------

\subsection{Online Learning}

Im Gegensatz zum \textbf{Batch Learning} (Trainingsdaten vollständig vorab
verfügbar) erhält der Lernalgorithmus beim \textbf{Online Learning} eine
Sequenz von Beispielen $(x^{(1)},y^{(1)}), (x^{(2)},y^{(2)}),\ldots$
nacheinander: für jedes $x^{(i)}$ wird zunächst eine Vorhersage getroffen,
bevor das wahre Label $y^{(i)}$ offengelegt wird (und ggf.\ für weiteres
Lernen genutzt werden kann). Von Interesse ist die \textbf{Gesamtzahl der
Fehler}, die der Algorithmus während dieses Prozesses macht.

\subsection{Der Perceptron-Algorithmus}

Mit Labels $y \in \{-1,+1\}$ und Parametern $\theta \in \R^{d+1}$ lautet die
Hypothese
\begin{equation}
  \htheta(x) = g(\theta\T x), \qquad
  g(z) = \begin{cases} 1 & z \geq 0 \\ -1 & z < 0. \end{cases}
\end{equation}
Bei korrekter Klassifikation bleibt $\theta$ unverändert; bei einem Fehler
($\htheta(x)\neq y$ für Beispiel $(x,y)$) erfolgt das Update
\begin{equation}
  \theta \leftarrow \theta + y\, x.
\end{equation}
Anders als bei LMS oder logistischer Regression modelliert der Perceptron
keine Wahrscheinlichkeit -- er ist ein rein diskriminatives
Entscheidungsverfahren, das nur bei Fehlern reagiert.

\subsection{Mistake-Bound-Theorem}

\begin{tcolorbox}[thmstyle]
\textbf{Satz (Block 1962, Novikoff 1962).}
Sei eine Sequenz $(x^{(1)},y^{(1)}),\ldots,(x^{(m)},y^{(m)})$ gegeben mit
$\norm{x^{(i)}} \leq D$ für alle $i$, und es existiere ein Einheitsvektor $u$
($\norm{u}=1$) mit $y^{(i)}(u\T x^{(i)}) \geq \gamma$ für alle $i$ (d.\,h.\
$u$ trennt die Daten mit Margin $\gamma$). Dann macht der
Perceptron-Algorithmus höchstens $(D/\gamma)^2$ Fehler auf dieser Sequenz.
\end{tcolorbox}

Bemerkenswert: Die Schranke hängt \emph{weder} von der Anzahl der Beispiele
$m$ noch von der Dimension $d$ ab -- nur vom (normierten) Margin $\gamma$ und
dem Datenradius $D$. \textbf{Beweisidee:} Sei $\theta^{(k)}$ der Parameter
beim $k$-ten Fehler ($\theta^{(1)}=\vec 0$). Einerseits wächst
$(\theta^{(k+1)})\T u \geq k\gamma$ (jedes Update erhöht die Projektion auf
$u$ um mind.\ $\gamma$), andererseits wächst $\norm{\theta^{(k+1)}}^2 \leq
kD^2$ (da bei einem Fehler $y^{(i)}(x^{(i)})\T\theta^{(k)} \leq 0$ gilt, folgt
$\norm{\theta^{(k+1)}}^2 \leq \norm{\theta^{(k)}}^2 + D^2$). Mit
Cauchy-Schwarz ($\norm{\theta^{(k+1)}}\geq(\theta^{(k+1)})\T u$, da $u$
Einheitsvektor) folgt $\sqrt{k}D \geq k\gamma$, also $k \leq (D/\gamma)^2$.

% ------------------------------------------------------------------
\section{Decision Trees}
% ------------------------------------------------------------------

\subsection{Nichtlineare, regionsbasierte Klassifikation}

Decision Trees sind die ersten genuin \textbf{nichtlinearen} Verfahren in
dieser Zusammenfassung: Sie partitionieren den Eingaberaum $\X$ rekursiv in
disjunkte Regionen, $\X = \bigcup_i R_i$ mit $R_i \cap R_j = \emptyset$ für
$i\neq j$, und sagen pro Region einen konstanten Wert voraus (Mehrheitsklasse
bei Klassifikation, Mittelwert bei Regression).

\subsection{Rekursive Partitionierung}

Da optimale Regionen im Allgemeinen intraktabel zu bestimmen sind, wählen
Decision Trees einen \textbf{gierigen, Top-Down}-Ansatz: Gegeben eine
Elternregion $R_p$, Feature-Index $j$ und Schwellwert $t$:
\begin{equation}
  R_1 = \{x \in R_p : x_j < t\}, \qquad R_2 = \{x \in R_p : x_j \geq t\}.
\end{equation}
Bei jedem Schritt werden Blattregion, Feature und Schwelle so gewählt, dass
der \emph{Rückgang} des kardinalitätsgewichteten Verlusts maximal wird:
\begin{equation}
  \max_{R_p,\,j,\,t} \; L(R_p) - \frac{|R_1|\,L(R_1) + |R_2|\,L(R_2)}{|R_1|+|R_2|}.
\end{equation}

\subsection{Verlustfunktionen für Splits}

Für Klassifikation mit Klassenanteil $\hat p_c$ (Anteil der Beispiele der
Klasse $c$) in Region $R$:
\begin{align}
  \text{Missklassifikationsverlust:} &\quad L_\text{misclass}(R) = 1 - \max_c \hat p_c, \\
  \text{Cross-Entropy:} &\quad L_\text{cross}(R) = -\sum_c \hat p_c \log_2 \hat p_c.
\end{align}

\begin{tcolorbox}[defstyle]
\textbf{Warum Cross-Entropy statt Missklassifikationsverlust?} Der
Missklassifikationsverlust ist nicht \emph{strikt konkav} -- ein Split kann
daher den gewichteten Verlust exakt unverändert lassen, obwohl er informativ
ist (Beispiel: Elternregion $400{+}/100{-}$; Split A in $150{+}/100{-}$ und
$250{+}/0{-}$ hat denselben gewichteten Missklassifikationsverlust wie ein
uninformativer Split B). Cross-Entropy (und Gini-Loss) sind strikt konkav:
Solange $\hat p_1 \neq \hat p_2$ in nicht-leeren Kindern, sinkt der
gewichtete Verlust garantiert gegenüber dem Elternknoten. Der Rückgang heißt
\textbf{Information Gain}.
\end{tcolorbox}

Für Regression wird pro Blatt der Mittelwert
$\hat y = \tfrac{1}{|R|}\sum_{i\in R} y_i$ vorhergesagt, und der quadratische
Verlust $L_\text{squared}(R) = \tfrac{1}{|R|}\sum_{i \in R}(y_i-\hat y)^2$
steuert die Split-Wahl.

\subsection{Regularisierung und Eigenschaften}

Ein voll ausgewachsener Baum (ein Trainingspunkt pro Blatt) hat niedrigen
Bias, aber hohe Varianz. Übliche Stop- und Pruning-Heuristiken:
\begin{itemize}
  \item Mindestgröße pro Blatt, maximale Tiefe, maximale Blattzahl.
  \item \textbf{Pruning}: Baum zunächst voll wachsen lassen, dann Knoten
        entfernen, die den Validierungsfehler minimal erhöhen -- robuster
        als ein Mindest-Loss-Rückgang pro Split, da sonst Interaktionseffekte
        höherer Ordnung (die erst nach mehreren Splits sichtbar werden)
        übersehen würden.
\end{itemize}
Laufzeit: $O(d)$ zur Vorhersage (Baumtiefe $d$), $O(nfd)$ zum Training
($n$ Beispiele, $f$ Features). \textbf{Vorteile}: interpretierbar,
verarbeitet kategoriale Features nativ (Splits über Teilmengen), schnell.
\textbf{Nachteile}: hohe Varianz, kann additive Struktur (z.\,B.\ eine
Grenze der Form $x_1+x_2=c$) nur durch viele Splits approximieren, obwohl
ein lineares Modell sie direkt abbilden könnte.

% ------------------------------------------------------------------
\section{Ensemble-Methoden}
% ------------------------------------------------------------------

\subsection{Varianzreduktion durch Mittelung}

Für $n$ identisch verteilte, paarweise korrelierte Zufallsvariablen (z.\,B.\
die Fehler einzelner Modelle) mit Varianz $\sigma^2$ und Korrelation $\rho$
gilt für den Mittelwert:
\begin{equation}
  \operatorname{Var}(\bar X) = \rho\sigma^2 + \frac{1-\rho}{n}\sigma^2.
\end{equation}
Mehr Modelle ($n\uparrow$) lässt den zweiten Term verschwinden; geringere
Korrelation zwischen den Modellen ($\rho\downarrow$) senkt beide Terme.
Ensembling zielt daher darauf ab, \textbf{viele, wenig korrelierte} Modelle
zu kombinieren.

\subsection{Bagging (Bootstrap Aggregation)}

\textbf{Bootstrap}: Ziehe aus dem Trainingsset $S$ ($|S|=n$) wiederholt $M$
Stichproben $Z_1,\ldots,Z_M$ \emph{mit Zurücklegen} (je Größe $n$). Trainiere
je ein Modell $G_m$ auf $Z_m$ und mittle:
\begin{equation}
  G(x) = \frac{1}{M}\sum_{m=1}^M G_m(x).
\end{equation}
Bagging senkt $\rho$ (Modelle sehen unterschiedliche Stichproben) und damit
die Varianz; der Bias steigt leicht (jede Bootstrap-Stichprobe enthält im
Schnitt nur $\approx 63\,\%$ der eindeutigen Trainingspunkte), was in der
Praxis meist durch den Varianzgewinn aufgewogen wird. Mehr Modelle können
\emph{nicht} zu mehr Overfitting führen, da $\rho$ unabhängig von $M$ ist.
\textbf{Out-of-Bag-Schätzung}: Das im Schnitt $\approx 1/3$ nicht gezogene
Restdrittel jeder Stichprobe dient als kostenlose Validierungsmenge -- im
Grenzwert $M\to\infty$ äquivalent zu Leave-One-Out Cross-Validation.

\begin{tcolorbox}[defstyle]
\textbf{Random Forests.} Bagging von Decision Trees, wobei an jedem Split
zusätzlich nur eine \emph{zufällige Teilmenge} der Features zur Auswahl
steht. Dies dekorreliert die Bäume weiter (ein einzelner starker Prädiktor
dominiert nicht mehr jeden Baum) und erlaubt robusteren Umgang mit fehlenden
Features, da nicht jeder Baum von jedem Feature abhängt.
\end{tcolorbox}

\subsection{Boosting}

Während Bagging Varianz reduziert, reduziert \textbf{Boosting} den
\textbf{Bias}: Es kombiniert viele \textbf{schwache Lerner} (hoher Bias,
niedrige Varianz -- z.\,B.\ Decision Stumps, Bäume der Tiefe~1) sequentiell,
wobei jeder neue Lerner sich verstärkt auf die Fehler der bisherigen
Ensemble-Vorhersage konzentriert.

\begin{tcolorbox}[algostyle, title={AdaBoost}]
\begin{enumerate}[leftmargin=*, itemsep=2pt]
  \item Initialisiere Gewichte $w_i \leftarrow 1/N$ für alle $i=1,\ldots,N$.
  \item Für $m = 1,\ldots,M$:
  \begin{enumerate}[label=(\alph*)]
    \item Trainiere schwachen Klassifikator $G_m$ auf den gewichteten Daten.
    \item Gewichteter Fehler: $\text{err}_m = \dfrac{\sum_i w_i \I\{y_i \neq G_m(x_i)\}}{\sum_i w_i}$.
    \item Klassifikatorgewicht: $\alpha_m = \log\!\bigl(\tfrac{1-\text{err}_m}{\text{err}_m}\bigr)$.
    \item Update: $w_i \leftarrow w_i \exp\bigl(\alpha_m \I\{y_i\neq G_m(x_i)\}\bigr)$.
  \end{enumerate}
  \item Ausgabe: $f(x) = \operatorname{sign}\bigl(\sum_{m=1}^M \alpha_m G_m(x)\bigr)$.
\end{enumerate}
\end{tcolorbox}

Falsch klassifizierte Beispiele erhalten höheres Gewicht -- nachfolgende
Lerner konzentrieren sich auf die "schwierigen" Fälle. Da spätere Lerner von
früheren abhängen (im Gegensatz zur Unabhängigkeit bei Bagging), steigt mit
wachsendem $M$ das Overfitting-Risiko.

\subsection{Forward Stagewise Additive Modeling}

AdaBoost ist ein Spezialfall des allgemeineren Schemas: Baue
$f_m(x) = f_{m-1}(x) + \beta_m G(x;\gamma_m)$ iterativ auf, wobei bei jedem
Schritt
\begin{equation}
  (\beta_m,\gamma_m) = \arg\min_{\beta,\gamma} \sum_{i=1}^N
    L\bigl(y_i,\, f_{m-1}(x_i) + \beta\,G(x_i;\gamma)\bigr)
\end{equation}
gelöst wird. Für quadratischen Verlust entspricht dies dem Fitten auf die
Residuen $y_i - f_{m-1}(x_i)$; für binäre Klassifikation mit
\textbf{exponentiellem Verlust} $L(y,\hat y) = \exp(-y\hat y)$ ergibt sich
exakt AdaBoost. \textbf{Gradient Boosting} verallgemeinert dies auf beliebige
differenzierbare Verluste: Bei jedem Schritt wird ein neuer schwacher Lerner
an den negativen Gradienten $g_i = -\partial L(y_i,f(x_i))/\partial f(x_i)$
angepasst (Regression von $G(x_i;\gamma)$ auf $g_i$ per Least Squares).

\begin{tcolorbox}[webstyle]
Moderne Gradient-Boosting-Implementierungen -- \textbf{XGBoost} (Chen \&
Guestrin, 2016), \textbf{LightGBM} und \textbf{CatBoost} -- dominieren bis
heute Kaggle-Wettbewerbe und produktive ML-Systeme auf strukturierten/
tabellarischen Daten (Betrugserkennung, Kreditrisiko, Nachfrageprognosen);
laut mehreren Quellen stecken sie hinter der Mehrzahl der Gewinnerlösungen
für tabellarische Aufgaben. Sie erweitern das hier vorgestellte
Gradient-Boosting-Schema um explizite Regularisierung, effizientes
histogrammbasiertes Splitting sowie native Unterstützung fehlender bzw.\
kategorialer Werte. Random Forests (Breiman, 2001) selbst bauen auf Bagging
(Breiman, 1996) und zufälliger Feature-Auswahl bei Splits (Ho, 1995/1998;
Amit \& Geman, 1997) auf.
\end{tcolorbox}

% ------------------------------------------------------------------
\section{Verlustfunktionen: Konvexe Ersatzverluste}
% ------------------------------------------------------------------

\subsection{Margin-basierte Verlustfunktionen}

Für binäre Klassifikation mit $y \in \{-1,+1\}$ heißt $z = y\,\theta\T x$ der
\textbf{Margin} eines Beispiels: $z>0$ bedeutet korrekte Klassifikation. Der
naheliegende Zielverlust wäre der \textbf{0/1-Verlust}
$\phi_{01}(z) = \I\{z\leq 0\}$ -- er ist jedoch unstetig, nicht konvex und
NP-hart zu minimieren. Man verwendet daher \textbf{konvexe Ersatzverluste}
(surrogate losses) $\phi(z)$ mit $\phi(z)\to 0$ für $z\to\infty$ und
$\phi(z)\to\infty$ für $z\to -\infty$:

\begin{center}
\begin{tabular}{lll}
\toprule
Verlust & $\phi(z)$ & Zugehörige Methode \\
\midrule
Logistisch & $\log(1+e^{-z})$ & Logistische Regression (Abschnitt~2) \\
Hinge & $\max(1-z,\,0)$ & Support Vector Machine (Abschnitt~6) \\
Exponentiell & $e^{-z}$ & AdaBoost (Abschnitt~9) \\
\bottomrule
\end{tabular}
\end{center}

\begin{tcolorbox}[thmstyle]
\textbf{Verbindendes Prinzip.} Logistische Regression, SVMs und Boosting
lösen scheinbar verschiedene Probleme -- tatsächlich minimieren alle drei
Verfahren $\frac{1}{n}\sum_i \phi\bigl(y^{(i)}\theta\T x^{(i)}\bigr)$ für
dieselbe Größe (den Margin), nur mit unterschiedlichem $\phi$. Alle drei
Verluste sind konvexe obere Schranken an den 0/1-Verlust und daher effizient
optimierbar, während sie dessen Zielrichtung (Vorzeichen des Margins)
approximieren.
\end{tcolorbox}

% ==================================================================
\newpage
\part{Deep Learning}
% ==================================================================

% ------------------------------------------------------------------
\section{Deep Learning}
% ------------------------------------------------------------------

\subsection{Nichtlineare Modelle und neuronale Netze}

\textbf{Lineare Modelle haben beschränkte Ausdruckskraft.} Durch Komposition
nichtlinearer Schichten kann ein neuronales Netz beliebig komplexe Funktionen
approximieren. Ein \textbf{Fully-Connected Netz} mit $r$ Schichten ist
rekursiv definiert:
\begin{align}
  a^{[0]} &= x, \\
  z^{[\ell]} &= W^{[\ell]}\, a^{[\ell-1]} + b^{[\ell]}, \\
  a^{[\ell]} &= \sigma\!\bigl(z^{[\ell]}\bigr), \quad \ell = 1,\ldots,r-1,\\
  \hat y &= W^{[r]}\, a^{[r-1]} + b^{[r]}.
\end{align}

Ohne nichtlineare Aktivierungsfunktion $\sigma$ kollabiert das Netz zu einem
linearen Modell: $W^{[r]}\cdots W^{[1]}x = \tilde W x$.

\subsection{Aktivierungsfunktionen}

\begin{center}
\begin{tabular}{lll}
\toprule
Name & Formel & Typischer Einsatz \\
\midrule
Sigmoid & $\sigma(z) = 1/(1+e^{-z})$ & Ausgabeschicht (binäre Klassifikation) \\
Tanh & $(e^z - e^{-z})/(e^z + e^{-z})$ & Veraltete versteckte Schichten \\
ReLU & $\max(z, 0)$ & Versteckte Schichten (Standard) \\
Leaky ReLU & $\max(z, \gamma z)$, $\gamma \in (0,1)$ & Verhindert ``Dying ReLU'' \\
GELU & $z \cdot \Phi(z)$ & Transformer-Modelle \\
Softmax & $e^{z_k}/\sum_j e^{z_j}$ & Ausgabeschicht (Mehrklassen) \\
\bottomrule
\end{tabular}
\end{center}

\subsection{Moderne Architekturbausteine}

\subsubsection*{Residualverbindungen (ResNets)}

Um den \textbf{Vanishing-Gradient}-Effekt in tiefen Netzen zu bekämpfen,
addiert ResNet den Eingabewert direkt zur Ausgabe eines Blocks:
\begin{equation}
  \text{Res}(z) = z + \sigma\!\bigl(W_2\,\sigma(W_1 z + b_1) + b_2\bigr).
\end{equation}
Der Gradient kann so ungehindert durch viele Schichten fließen.

\subsubsection*{Layer-Normalisierung}

\begin{equation}
  \mathrm{LN}(z) = \gamma \odot \frac{z - \hat\mu}{\hat\sigma} + \beta,
\end{equation}
wobei $\hat\mu, \hat\sigma$ Mittelwert und Standardabweichung der Einträge
von $z$ sind. Wichtige Eigenschaft: Skaleninvarianz bezüglich der
Gewichtsmatrizen.

\subsubsection*{Faltungsschichten (CNNs)}

Eine 1D-Faltungsschicht mit Filter $w \in \R^k$ und Stride $s$ berechnet:
\begin{equation}
  \mathrm{Conv}(z)_i = \sum_{j=1}^k w_j \, z_{i \cdot s - \ell + (j-1)}.
\end{equation}
Parameter-Sharing reduziert die Parameterzahl von $O(m^2)$ auf $O(km)$.

\subsection{Backpropagation}

\begin{tcolorbox}[thmstyle]
\textbf{Theorem (Backpropagation).}
Lässt sich ein Rechenausdruck der Größe $N$ für den Verlust in einem Forward
Pass auswerten, so lässt sich der vollständige Gradient $\nabla_\theta J$
ebenfalls in $O(N)$ Zeit berechnen.
\end{tcolorbox}

\textbf{Algorithmisch.} Sei der Forward Pass:
\[
  u^{[0]} = x, \quad u^{[\ell]} = M_\ell(u^{[\ell-1]};\,\theta^{[\ell]}),
  \quad J = u^{[L]}.
\]
Der Backward Pass berechnet rückwärts:
\begin{align}
  \frac{\partial J}{\partial u^{[L]}} &= 1, \\
  \frac{\partial J}{\partial u^{[\ell-1]}}
    &= \underbrace{B\bigl[M_\ell,\, u^{[\ell-1]}\bigr]}_{\text{Backward-Funktion}}
       \!\Bigl(\frac{\partial J}{\partial u^{[\ell]}}\Bigr), \\
  \frac{\partial J}{\partial \theta^{[\ell]}}
    &= B_\theta\bigl[M_\ell,\, \theta^{[\ell]}\bigr]
       \!\Bigl(\frac{\partial J}{\partial u^{[\ell]}}\Bigr).
\end{align}

\subsection{Optimierungsverfahren}

\begin{tcolorbox}[algostyle, title={Mini-Batch SGD}]
\begin{enumerate}[leftmargin=*, itemsep=1pt]
  \item Initialisiere $\theta$ (z.\,B. zufällig)
  \item Wiederhole bis Konvergenz:
    \begin{enumerate}[label=(\alph*)]
      \item Sample Batch $\mathcal{B} \subset \{1,\ldots,n\}$, $|\mathcal{B}| = B$
      \item $g \leftarrow \frac{1}{B}\sum_{i \in \mathcal{B}} \nabla_\theta J^{(i)}(\theta)$
      \item $\theta \leftarrow \theta - \alpha\, g$
    \end{enumerate}
\end{enumerate}
\end{tcolorbox}

% ==================================================================
\newpage
\part{Generalisierung und Regularisierung}
% ==================================================================

% ------------------------------------------------------------------
\section{Generalisierung}
% ------------------------------------------------------------------

\subsection{Trainings- und Testfehler}

Das eigentliche Ziel beim maschinellen Lernen ist die Minimierung des
\textbf{Testfehlers} (auch: generalization error):
\begin{equation}
  L(\theta) = \E_{(x,y) \sim \mathcal{D}}\bigl[(y - h_\theta(x))^2\bigr].
\end{equation}
Der \textbf{Trainingsfehler} $J(\theta) = \frac{1}{n}\sum_i (y^{(i)} -
h_\theta(x^{(i)}))^2$ ist lediglich eine Approximation.

\begin{itemize}
  \item \textbf{Overfitting:} $J(\theta) \ll L(\theta)$ -- Modell zu komplex.
  \item \textbf{Underfitting:} $J(\theta) \approx L(\theta)$ groß -- Modell
        zu einfach.
\end{itemize}

\subsection{Bias-Varianz-Zerlegung}

\subsubsection*{Intuition: Was bedeuten Bias und Varianz?}

Betrachtet man denselben Datensatz $\{(x^{(i)},y^{(i)})\}$ und passt sowohl
ein lineares als auch ein Grad-5-Polynom an, zeigt sich ein charakteristischer
Unterschied im Fehlverhalten: Das lineare Modell liegt \emph{systematisch}
daneben -- selbst mit beliebig viel Trainingsdaten würde eine Gerade die
zugrunde liegende (nichtlineare) Struktur nicht erfassen können. Das
Grad-5-Polynom hingegen trifft die Trainingspunkte fast perfekt, aber die
konkrete Kurvenform hängt stark davon ab, \emph{welche} zufällige
Stichprobe man gerade gezogen hat -- ein neuer, gleich großer Datensatz aus
derselben Verteilung würde zu einer sichtbar anderen Kurve führen.

\begin{itemize}
  \item \textbf{Bias} = der Fehler, der \emph{übrig bliebe, selbst wenn man
        unendlich viele Trainingsdaten hätte}. Er entsteht rein dadurch, dass
        die gewählte Hypothesenklasse (z.\,B.\ "alle Geraden") die wahre
        Funktion $h^*$ grundsätzlich nicht abbilden kann -- ein
        Ausdrucksstärke-Problem, kein Datenmengenproblem. Mehr Daten helfen
        hier \emph{nicht}.
  \item \textbf{Varianz} = wie stark das gelernte Modell $h_S$ schwankt, wenn
        man den Trainingsdatensatz $S$ (fester Größe) wiederholt neu zieht.
        Sie entsteht dadurch, dass ein zu flexibles Modell "zufällige"
        Muster im jeweiligen Rauschen $\xi^{(i)}$ mitlernt, statt nur das
        eigentliche Signal. Mehr Daten \emph{reduzieren} die Varianz, da
        einzelne Rauschausschläge im Mittel weniger Gewicht bekommen.
\end{itemize}

\subsubsection*{Im Parameterraum betrachtet}

Dieselbe Intuition lässt sich direkt im Parameterraum $\theta$ visualisieren.
Sei $\theta^*$ der Parameter des \emph{bestmöglichen} Modells innerhalb der
gewählten Hypothesenklasse (also der Grenzwert von $\hat\theta_S$ bei
unendlich großem Trainingsset). Zieht man stattdessen $k$ verschiedene,
endlich große Trainingssets $S_1,\ldots,S_k$ derselben Größe und trainiert
jeweils separat, erhält man $k$ Schätzungen $\hat\theta_{S_1},\ldots,
\hat\theta_{S_k}$ -- eine "Punktwolke" um $\theta^*$:

\begin{center}
\begin{tabular}{cc}
\begin{tikzpicture}[scale=1.05]
  \draw[gray!35] (-1,-1) rectangle (1,1);
  \foreach \x/\y in {0.05/0.03,-0.08/0.06,0.03/-0.09,-0.06/-0.04,
                     0.09/-0.02,-0.03/0.08,0.06/-0.07,-0.05/-0.06}{
    \fill[red!70!black] (\x,\y) circle (0.045);
  }
  \fill[mainblue] (0,0) circle (0.05);
  \node[mainblue, right] at (0.1,0.12) {\scriptsize $\theta^*$};
  \node[below] at (0,-1.15) {\footnotesize niedriger Bias, niedrige Varianz};
\end{tikzpicture}
&
\begin{tikzpicture}[scale=1.05]
  \draw[gray!35] (-1,-1) rectangle (1,1);
  \foreach \x/\y in {0.4/0.3,-0.6/0.5,0.3/-0.7,-0.5/-0.3,
                     0.75/-0.2,-0.3/0.65,0.5/-0.5,-0.4/-0.55}{
    \fill[red!70!black] (\x,\y) circle (0.045);
  }
  \fill[mainblue] (0,0) circle (0.05);
  \node[mainblue, right] at (0.1,0.12) {\scriptsize $\theta^*$};
  \node[below] at (0,-1.15) {\footnotesize niedriger Bias, hohe Varianz};
\end{tikzpicture}
\\[1.4cm]
\begin{tikzpicture}[scale=1.05]
  \draw[gray!35] (-1,-1) rectangle (1,1);
  \foreach \x/\y in {0.55/0.43,0.42/0.46,0.5/0.31,0.44/0.34,
                     0.53/0.38,0.4/0.4,0.46/0.45,0.5/0.35}{
    \fill[red!70!black] (\x,\y) circle (0.045);
  }
  \fill[mainblue] (0,0) circle (0.05);
  \node[mainblue, left] at (-0.05,-0.12) {\scriptsize $\theta^*$};
  \node[below] at (0,-1.15) {\footnotesize hoher Bias, niedrige Varianz};
\end{tikzpicture}
&
\begin{tikzpicture}[scale=1.05]
  \draw[gray!35] (-1,-1) rectangle (1,1);
  \foreach \x/\y in {0.6/0.7,0.3/0.9,0.8/0.3,0.15/0.5,
                     0.85/0.6,0.35/0.15,0.7/0.85,0.25/0.3}{
    \fill[red!70!black] (\x,\y) circle (0.045);
  }
  \fill[mainblue] (0,0) circle (0.05);
  \node[mainblue, left] at (-0.05,-0.12) {\scriptsize $\theta^*$};
  \node[below] at (0,-1.15) {\footnotesize hoher Bias, hohe Varianz};
\end{tikzpicture}
\end{tabular}
\end{center}

\begin{tcolorbox}[defstyle]
Jeder rote Punkt ist die Parameterschätzung $\hat\theta_{S_i}$ aus
\emph{einem} zufällig gezogenen Trainingsdatensatz; der blaue Punkt markiert
$\theta^*$. \textbf{Bias} ist der Abstand zwischen $\theta^*$ und dem
\emph{Zentrum} der Punktwolke ($\E_S[\hat\theta_S]$) -- eine systematische
Verschiebung, die auch bei unendlich vielen Wiederholungen bestehen bliebe.
\textbf{Varianz} ist die \emph{Streuung} der Punktwolke um ihr eigenes
Zentrum -- unabhängig davon, ob dieses Zentrum überhaupt nahe an $\theta^*$
liegt.
\end{tcolorbox}

Die Verbindung zur Modellkomplexität: Ein \textbf{einfaches Modell} (z.\,B.\
lineare Regression) hat einen kleinen, stark eingeschränkten Parameterraum
-- unterschiedliche Trainingssets führen daher zu sehr ähnlichen Schätzungen
$\hat\theta_S$ (\textbf{niedrige Varianz}), aber falls die wahre Funktion
außerhalb dieses eingeschränkten Raums liegt, sitzt auch das Zentrum der
Punktwolke systematisch daneben (\textbf{hoher Bias}) -- typischerweise der
untere linke Fall oben. Ein \textbf{komplexes Modell} (z.\,B.\ Grad-5-
Polynom) hat einen großen, flexiblen Parameterraum, der $\theta^*$ im Mittel
gut treffen kann (\textbf{niedriger Bias}), aber jede einzelne Schätzung
"jagt" das Rauschen der jeweiligen Stichprobe und landet daher an sehr
unterschiedlichen Stellen (\textbf{hohe Varianz}) -- der obere rechte Fall.

\subsubsection*{Formale Zerlegung}

Für das Regressionsproblem lässt sich der erwartete MSE an einem
Punkt $x$ zerlegen:
\begin{equation}
  \E\bigl[(y - h_S(x))^2\bigr]
  = \underbrace{\sigma^2}_{\text{irreduzibles Rauschen}}
  + \underbrace{\bigl(h^*(x) - \bar h(x)\bigr)^2}_{\text{Bias}^2}
  + \underbrace{\E\bigl[(\bar h(x) - h_S(x))^2\bigr]}_{\text{Varianz}},
\end{equation}
wobei $h^*$ die wahre Funktion, $\bar h = \E_S[h_S]$ der Erwartungswert des
gelernten Modells über zufällige Trainingsdatensätze $S$ ist -- im
Parameterraum-Bild oben entspricht $\bar h$ also gerade dem Zentrum der
Punktwolke.

\begin{itemize}
  \item Einfache Modelle ($h$ wenig ausdrucksstark) $\Rightarrow$ hoher Bias.
  \item Komplexe Modelle (viele Parameter) $\Rightarrow$ hohe Varianz.
  \item Optimale Modellkomplexität balanciert beides.
\end{itemize}

\subsection{Double Descent}

Das klassische Bild (Testfehler fällt, dann steigt er bei Überanpassung) gilt
nur im \emph{classical regime} ($n \gg d$). Im modernen Deep-Learning-Regime
$d \gg n$ tritt ein zweites Absinken auf:

\begin{center}
\begin{tikzpicture}
\begin{axis}[
  width=12cm, height=5.5cm,
  xlabel={Modellkomplexität ($d$)},
  ylabel={Testfehler},
  xmin=0, xmax=10, ymin=0, ymax=3,
  xtick={5}, xticklabels={$n \approx d$},
  ytick=\empty,
  grid=both, grid style={line width=0.3pt, draw=gray!30},
  axis lines=left, tick label style={font=\small},
  label style={font=\small},
]
  \addplot[mainblue, very thick, domain=0:4.8, samples=80]
    {0.8 + 0.5*exp(-x) + 0.3*x^0.5};
  \addplot[mainblue, very thick, domain=4.8:5.1, samples=10]
    {0.8 + 0.5*exp(-x) + 0.3*x^0.5 + 8*(x-4.8)};
  \addplot[mainblue, very thick, domain=5.1:10, samples=80]
    {3.5*exp(-0.6*(x-5.1)) + 0.5};
  \draw[dotted, gray!70] (axis cs:5,0) -- (axis cs:5,3);
  \node[font=\scriptsize, above] at (axis cs:2.5, 1.4) {klassisch};
  \node[font=\scriptsize, above] at (axis cs:7.5, 1.2) {modern (überparametrisiert)};
\end{axis}
\end{tikzpicture}
\end{center}

Der Peak tritt auf, wenn $n \approx d$ -- das Modell interpoliert die
Trainingsdaten, verhält sich aber schlecht bei unregelmäßiger Interpolation.
Im überparametrisierten Regime wählt Gradient Descent implizit die
\emph{Minimum-Norm-Lösung}, die generell gut generalisiert.

\subsection{Konzentrationsungleichungen}

Grundlage der Stichprobenkomplexität ist die Frage: Wie nah liegt der
Trainingsfehler $J(\theta)$ am wahren Testfehler $L(\theta)$? Antwort liefern
\textbf{Konzentrationsungleichungen} -- eine Kette zunehmend schärferer
Abschätzungen (Markov $\to$ Chebyshev $\to$ Chernoff), deren wichtigster
Spezialfall für beschränkte Zufallsvariablen die Hoeffdingsche Ungleichung
ist.

\begin{tcolorbox}[thmstyle]
\textbf{Hoeffdingsche Ungleichung.} Seien $Z_1,\ldots,Z_n$ unabhängige,
beschränkte Zufallsvariablen mit $Z_i \in [a,b]$. Dann gilt für alle
$t\geq 0$:
\begin{equation}
  P\Bigl(\Bigl|\frac1n\sum_{i=1}^n (Z_i - \E[Z_i])\Bigr| \geq t\Bigr)
  \leq 2\exp\!\Bigl(-\frac{2nt^2}{(b-a)^2}\Bigr).
\end{equation}
\end{tcolorbox}

Für den Trainingsfehler einer \emph{festen} Hypothese $h$ mit
$Z_i = \I\{h(x^{(i)}) \neq y^{(i)}\} \in \{0,1\}$ liefert dies ($a=0,\,b=1$):
\begin{equation}
  P\bigl(|J(h) - L(h)| \geq t\bigr) \leq 2\exp(-2nt^2).
\end{equation}
Diese Schranke gilt jedoch nur für eine im Voraus fixierte Hypothese. Da
$\hat h$ erst \emph{nach} Betrachten der Trainingsdaten aus $\mathcal H$
ausgewählt wird, benötigt man eine \textbf{gleichmäßige Konvergenz}
("uniform convergence") über \emph{alle} $h\in\mathcal H$ simultan -- per
Union Bound über $|\mathcal H|=k$ Hypothesen ergibt sich (Auflösen von
$\delta = 2k\exp(-2nt^2)$ nach $t$) genau die Schranke der folgenden
Stichprobenkomplexität.

\subsection{Stichprobenkomplexität}

Für eine endliche Hypothesenklasse $\mathcal{H}$ mit $|\mathcal{H}| = k$
gilt mit Wahrscheinlichkeit $\geq 1-\delta$:
\begin{equation}
  L(\hat\theta) \leq J(\hat\theta) + \sqrt{\frac{\log(k/\delta)}{2n}}.
\end{equation}
Die Generalisierungslücke schrumpft mit $1/\sqrt{n}$ und wächst logarithmisch
mit $|\mathcal{H}|$. Für unendliche $\mathcal{H}$ (z.\,B.\ alle linearen
Klassifikatoren) übernimmt die \textbf{VC-Dimension} die Rolle von
$\log |\mathcal H|$.

\begin{tcolorbox}[defstyle]
\textbf{VC-Dimension (Vapnik--Chervonenkis).} Die VC-Dimension
$\mathrm{VC}(\mathcal H)$ einer Hypothesenklasse ist die Größe der größten
Punktmenge, die $\mathcal H$ in \emph{jeder} möglichen Beschriftung korrekt
klassifizieren kann ("shattern"). Für lineare Klassifikatoren in $\R^d$ gilt
$\mathrm{VC}(\mathcal H) = d+1$. Mit Wahrscheinlichkeit $\geq 1-\delta$ gilt
dann für \emph{alle} $h \in \mathcal H$ gleichzeitig:
\begin{equation}
  L(h) \leq J(h) + O\!\left(\sqrt{\frac{\mathrm{VC}(\mathcal H)}{n}
    \log\frac{n}{\mathrm{VC}(\mathcal H)} + \frac1n\log\frac1\delta}\right).
\end{equation}
Die für eine $\varepsilon$-genaue Schätzung benötigte Stichprobengröße wächst
linear mit $\mathrm{VC}(\mathcal H)$ -- bei linearen Modellen also linear mit
der Parameteranzahl $d$. Für Modellklassen ohne endliche VC-Dimension
(z.\,B.\ manche Kernel- oder Deep-Learning-Modelle) versagt diese klassische
Theorie -- ein Grund, warum das Double-Descent-Phänomen (vorheriger
Abschnitt) lange Zeit überraschend war.
\end{tcolorbox}

% ------------------------------------------------------------------
\section{Regularisierung und Modellauswahl}
% ------------------------------------------------------------------

\subsection{Regularisierung}

\textbf{Regularisierung} fügt der Verlustfunktion einen Strafterm hinzu, der
große Gewichte bestraft:
\begin{align}
  \text{L2 (Ridge):} &\quad J_\lambda(\theta) = J(\theta) + \tfrac{\lambda}{2}\norm{\theta}^2, \\
  \text{L1 (Lasso):} &\quad J_\lambda(\theta) = J(\theta) + \lambda\norm{\theta}_1.
\end{align}

L2-Regularisierung fördert gleichmäßig kleine Gewichte. L1 fördert
\textbf{Sparsität} (viele Gewichte exakt 0), weil die $|\theta_j|$-Strafe
eine Unstetigkeit im Subgradienten bei $\theta_j = 0$ erzeugt.

\subsection{Bayesianische Interpretation}

\begin{tcolorbox}[defstyle]
Aus Bayesianischer Sicht entspricht L2-Regularisierung einem
\textbf{Gauß-Prior} $\theta_j \sim \N(0, \tau^2)$:
\[
  \theta_\mathrm{MAP}
  = \arg\max_\theta \bigl[\log p(y \mid X,\theta) + \log p(\theta)\bigr]
  = \arg\min_\theta \bigl[\text{NLL}(\theta) + \tfrac{1}{2\tau^2}\norm{\theta}^2\bigr].
\]
L1-Regularisierung entspricht einem \textbf{Laplace-Prior}
$p(\theta_j) \propto e^{-|\theta_j|/b}$.
\end{tcolorbox}

Für lineare Regression mit Gauß-Prior gibt es eine geschlossene Lösung:
\begin{equation}
  \theta_\mathrm{MAP} = (X\T X + 2\lambda\sigma^2 I)^{-1} X\T \bm{y}.
\end{equation}
Der Regularisierungsterm $2\lambda\sigma^2 I$ stellt sicher, dass die Matrix
stets invertierbar ist.

\subsection{Implizite Regularisierung}

Auch ohne expliziten Regularisierungsterm wirkt Gradient Descent als
impliziter Regularisierer: Im überparametrisierten Regime (mehr Parameter
als Datenpunkte) wählt SGD unter allen Interpolatoren bevorzugt denjenigen
mit minimaler Norm -- was einer L2-Regularisierung ähnelt.

\subsection{Kreuzvalidierung (Cross-Validation)}

Zur Auswahl von Hyperparametern (z.\,B.\ $\lambda$, Lernrate, Netzgröße)
wird \textbf{K-Fold Cross-Validation} eingesetzt:
\begin{enumerate}
  \item Teile Trainingsdaten in $K$ gleich große Teile.
  \item Trainiere $K$-mal, je einmal auf $K-1$ Teilen.
  \item Validierungsfehler auf dem ausgelassenen Teil messen.
  \item Mittelwert über alle $K$ Durchläufe als Schätzer des Testfehlers.
\end{enumerate}
Wähle den Hyperparameter mit kleinstem mittleren Validierungsfehler, dann
trainiere auf dem \emph{gesamten} Trainingsdatensatz neu.

% ------------------------------------------------------------------
\section{Gaussian Processes}
% ------------------------------------------------------------------

\begin{tcolorbox}[srcstyle]
Quelle: \emph{Gaussian processes}, Chuong B. Do (aktualisiert von Honglak
Lee), CS229 Section Notes -- \texttt{section/cs229-gaussian\_processes.pdf}.
\end{tcolorbox}

\subsection{Bayessche Sichtweise auf Regression}

Klassische Verfahren (lineare Regression, SVM) liefern einen einzelnen
"Best-Fit"-Parametervektor $\theta$ bzw.\ eine Punktschätzung $h(x)$.
\textbf{Bayessche Methoden} liefern stattdessen eine
\textbf{Posterior-Verteilung} über Modelle -- und damit direkt ein Maß für
die Unsicherheit einer Vorhersage. \textbf{Gaussian-Process-Regression} ist
ein kernelbasiertes, voll bayessches Regressionsverfahren, das die
probabilistische lineare Regression (Abschnitt~1.4), die Bayesianische
Interpretation der Regularisierung (Abschnitt~13.2) und Kernels
(Abschnitt~5) zusammenführt.

\subsection{Verteilungen über Funktionen}

Ein \textbf{Gaussian Process (GP)} ist ein stochastischer Prozess
$\{f(x) : x\in\X\}$, sodass jede endliche Teilkollektion
$f(x^{(1)}),\ldots,f(x^{(m)})$ gemeinsam multivariat gaußverteilt ist. Er
wird vollständig durch eine Mittelwertfunktion $m(\cdot)$ und eine
Kovarianzfunktion -- genau den bekannten \textbf{Kernel} $k(\cdot,\cdot)$! --
festgelegt:
\begin{equation}
  f(\cdot) \sim \mathcal{GP}(m(\cdot),\,k(\cdot,\cdot)),
\end{equation}
sodass für beliebige Punkte $x^{(1)},\ldots,x^{(m)}$ gilt
\begin{equation}
  \begin{pmatrix} f(x^{(1)}) \\ \vdots \\ f(x^{(m)}) \end{pmatrix}
  \sim \N\!\left(
    \begin{pmatrix} m(x^{(1)}) \\ \vdots \\ m(x^{(m)}) \end{pmatrix},\;
    \begin{pmatrix}
      k(x^{(1)},x^{(1)}) & \cdots & k(x^{(1)},x^{(m)}) \\
      \vdots & \ddots & \vdots \\
      k(x^{(m)},x^{(1)}) & \cdots & k(x^{(m)},x^{(m)})
    \end{pmatrix}
  \right).
\end{equation}
Mit dem \textbf{Squared-Exponential-Kernel}
$k(x,x') = \exp\bigl(-\norm{x-x'}^2/(2\tau^2)\bigr)$ (identisch zum
Gauß-/RBF-Kernel aus Abschnitt~5.2) sind Funktionswerte nahe beieinander
liegender Punkte stark korreliert -- Stichproben aus dem GP-Prior sind
"lokal glatt", und die Korrelation fällt mit dem Abstand im Eingaberaum ab.

\subsection{GP-Regression: Vorhersage}

Modell: $y^{(i)} = f(x^{(i)}) + \varepsilon^{(i)}$ mit
$\varepsilon^{(i)} \overset{\text{iid}}{\sim} \N(0,\sigma^2)$ und
Nullmittel-GP-Prior $f(\cdot)\sim\mathcal{GP}(0,k(\cdot,\cdot))$. Für
Trainingsdaten $X$ (Ausgaben $\vec y$) und Testpunkte $X_*$ sind $\vec y$ und
die gesuchten Funktionswerte $\vec f_*$ am Testpunkt gemeinsam gaußverteilt:
\begin{equation}
  \begin{pmatrix}\vec y \\ \vec f_*\end{pmatrix}
  \sim \N\!\left(\vec 0,\;
    \begin{pmatrix} K(X,X)+\sigma^2 I & K(X,X_*) \\ K(X_*,X) & K(X_*,X_*)\end{pmatrix}
  \right),
\end{equation}
wobei $(K(X,X))_{ij} = k(x^{(i)},x^{(j)})$ die \textbf{Gram-Matrix} ist
(vgl.\ Mercers Satz, Abschnitt~5.3).

\begin{tcolorbox}[thmstyle]
\textbf{Posterior Predictive Distribution.} Durch Konditionierung der
gemeinsamen Gauß-Verteilung (Formeln aus dem Faktorenanalyse-Abschnitt)
ergibt sich in geschlossener Form:
\begin{align}
  \vec f_* \mid \vec y, X, X_* &\sim \N(\mu_*,\,\Sigma_*), \\
  \mu_* &= K(X_*,X)\bigl(K(X,X)+\sigma^2 I\bigr)^{-1}\vec y, \\
  \Sigma_* &= K(X_*,X_*) - K(X_*,X)\bigl(K(X,X)+\sigma^2 I\bigr)^{-1}K(X,X_*).
\end{align}
\end{tcolorbox}

Bemerkenswert: Trotz der scheinbaren Komplexität von "Verteilungen über
Funktionen" ist die Vorhersage nur eine Anwendung der
Konditionierungsformel für multivariate Gauß-Verteilungen. $\mu_*$ liefert
die Punktschätzung, $\Sigma_*$ ein kalibriertes Unsicherheitsmaß
(Konfidenzband), das klassische (nicht-bayessche) Kernel-Regression nicht
liefert -- der zentrale praktische Vorteil von GPs.

\begin{tcolorbox}[defstyle]
\textbf{Bezug zur Bayesianischen Regularisierung.} Bayessche lineare
Regression mit Gauß-Prior (Abschnitt~13.2) liefert, korrekt "kernelisiert",
exakt dieselbe Vorhersageverteilung wie GP-Regression -- die GP-Sichtweise
("Prior direkt über Funktionen statt über Parameter $\theta$") liefert dafür
aber eine deutlich einfachere Herleitung.
\end{tcolorbox}

% ------------------------------------------------------------------
\section{Praktische Fehleranalyse}
% ------------------------------------------------------------------

Theoretische Bias-Varianz-Analyse beantwortet \emph{ob} ein Modell zu einfach
oder zu komplex ist -- in der Praxis muss man bei mehrstufigen Systemen oft
zusätzlich herausfinden, \emph{welche Komponente} für einen hohen Fehler
verantwortlich ist, um die eigene Arbeitszeit sinnvoll zu priorisieren.

\subsection{Fehleranalyse (Error Analysis)}

Bei einer Pipeline aus mehreren Komponenten (Beispiel Gesichtserkennung:
Preprocessing $\to$ Gesichtsdetektion $\to$ Augen-/Nase-/Mund-Erkennung
$\to$ finale Klassifikation) ersetzt man nacheinander \emph{jede} Komponente
durch ihre \textbf{Ground-Truth-Ausgabe} (von Hand annotiert) und misst die
resultierende Gesamtgenauigkeit:

\begin{center}
\begin{tabular}{lc}
\toprule
Komponente mit perfekter (Ground-Truth-)Ausgabe & Genauigkeit \\
\midrule
Ausgangssystem (keine Komponente ersetzt) & 85\,\% \\
+ perfektes Preprocessing (Hintergrundentfernung) & 85.1\,\% \\
+ perfekte Gesichtsdetektion & 91\,\% \\
+ perfekte Augen-Segmentierung & 95\,\% \\
\bottomrule
\end{tabular}
\end{center}

Der Sprung bei einer Komponente zeigt, wie viel \emph{maximal} durch deren
Verbesserung zu gewinnen ist -- ein kleiner Sprung (Preprocessing: nur
$+0.1$ Prozentpunkte) signalisiert, dass sich weitere Arbeit an dieser
Komponente kaum lohnt, selbst wenn sie perfektioniert würde.

\subsection{Ablative Analyse}

Fehleranalyse geht vom Ausgangssystem zum perfekten System; \textbf{ablative
Analyse} geht umgekehrt vom \emph{aktuellen, bereits guten} System aus und
entfernt schrittweise Komponenten (z.\,B.\ Features eines Spam-Klassifikators:
Rechtschreibkorrektur, Sender-Host-Features, Header-Features, \ldots), um zu
zeigen, wie viel jede einzelne Komponente zur finalen Performance beigetragen
hat. Dies ist besonders nützlich, um in einer Publikation zu begründen,
welche Modellbestandteile tatsächlich wichtig waren, statt nur die
Gesamtgenauigkeit anzugeben.

\subsection{Fehler manuell analysieren}

Ergänzend: eine zufällige Stichprobe (z.\,B.\ 100) fehlklassifizierter
Beispiele manuell inspizieren und nach Fehlerkategorien gruppieren (z.\,B.\
"unscharfe Bilder", "Verwechslung mit ähnlicher Klasse"). Macht eine
Kategorie $x\,\%$ der beobachteten Fehler aus, ist $x\,\%$ die \emph{obere
Schranke} für den durch Behebung dieser Kategorie erzielbaren
Genauigkeitsgewinn -- so lässt sich die eigene Arbeitszeit nach erwartetem
Ertrag statt nach Bauchgefühl priorisieren.

\subsection{Bias-Varianz-Diagnostik via Lernkurven}

Statt planlos "irgendetwas" am Modell zu ändern (mehr Daten? andere
Features? SVM statt logistischer Regression?), empfiehlt sich zunächst
eine gezielte \textbf{Diagnose}: Man plottet Trainings- und Testfehler als
Funktion der Trainingsgröße $m$ (\textbf{Lernkurve}):

\begin{itemize}
  \item \textbf{Hohe Varianz}: großer Abstand zwischen Trainings- und
        Testfehler; Testfehler fällt mit wachsendem $m$ noch deutlich --
        mehr Trainingsdaten helfen voraussichtlich.
  \item \textbf{Hoher Bias}: bereits der Trainingsfehler ist inakzeptabel
        hoch, der Abstand zum Testfehler ist klein -- mehr Daten helfen
        \emph{nicht}, es braucht ein ausdrucksstärkeres Modell (mehr
        Features, komplexere Hypothesenklasse).
\end{itemize}

\begin{tcolorbox}[defstyle]
\textbf{Welche Maßnahme behebt was?} Mehr Trainingsbeispiele und ein
kleinerer Feature-Satz reduzieren \emph{Varianz}; ein größerer Feature-Satz
oder ein ausdrucksstärkeres Modell reduziert \emph{Bias}. Diese Zuordnung
lässt sich aus der Lernkurve direkt ablesen, bevor man Zeit in die
Umsetzung investiert.
\end{tcolorbox}

\subsection{Diagnostik: Optimierungsalgorithmus vs.\ Optimierungsziel}

Ein zweites, subtileres Diagnose-Muster: Angenommen ein SVM erzielt bessere
Ergebnisse als bayessche logistische Regression (BLR) bezüglich der
eigentlich interessierenden Metrik $a(\theta)$ (z.\,B.\ gewichtete
Genauigkeit) -- obwohl eigentlich BLR eingesetzt werden soll (z.\,B.\ aus
Effizienzgründen). BLR optimiert dabei intern nicht $a(\theta)$, sondern die
Log-Likelihood $J(\theta)$. Um zu verstehen, \emph{woran} es liegt, vergleicht
man $J(\theta_\text{SVM})$ mit $J(\theta_\text{BLR})$:

\begin{itemize}
  \item $J(\theta_\text{BLR}) < J(\theta_\text{SVM})$: BLR hat sein
        \emph{eigenes} Ziel $J$ nicht einmal erreicht $\Rightarrow$ Problem
        liegt beim \textbf{Optimierungsalgorithmus} (Gradientenverfahren
        nicht konvergiert). Abhilfe: mehr Iterationen, Newton-Verfahren,
        anderer Lernrate.
  \item $J(\theta_\text{BLR}) \geq J(\theta_\text{SVM})$: BLR hat $J$
        erfolgreich maximiert, trotzdem schneidet es bei $a(\theta)$
        schlechter ab als die SVM $\Rightarrow$ Problem liegt beim
        \textbf{Optimierungsziel} selbst: $J$ ist nicht die richtige
        Zielfunktion für das, was man eigentlich erreichen will. Abhilfe:
        anderes $\lambda$/$C$, andere Modellklasse, gewichtete
        Verlustfunktion.
\end{itemize}

Diese Unterscheidung verhindert, dass man z.\,B.\ stundenlang an der
Konvergenz eines Gradientenverfahrens optimiert, obwohl in Wahrheit die
Zielfunktion selbst falsch gewählt war (oder umgekehrt).

\begin{tcolorbox}[srcstyle]
Quelle: \emph{Advice for applying Machine Learning}, Andrew Ng (CS229
Vorlesungsfolien) -- \texttt{materials/ML-advice.pdf}.
\end{tcolorbox}

% ------------------------------------------------------------------
\section{Evaluationsmetriken für Klassifikation}
% ------------------------------------------------------------------

\begin{tcolorbox}[srcstyle]
Quelle: \emph{Evaluation Metrics}, Yining Chen (adaptiert von Folien von
Anand Avati), CS229 Section Notes --
\texttt{section/evaluation\_metrics\_spring2020.pdf}.
\end{tcolorbox}

\subsection{Warum Genauigkeit oft nicht genügt}

Die Trainings-Kostenfunktion ist nur ein \emph{Proxy} für das eigentliche
Geschäfts-/Anwendungsziel. Evaluationsmetriken übersetzen dieses Ziel in
eine quantitative, vergleichbare Größe -- besonders wichtig, wenn nicht
alle Fehlerarten gleich teuer sind (z.\,B.\ ein übersehener Krebsbefund vs.\
ein falscher Alarm).

\subsection{Konfusionsmatrix und Punkt-Metriken}

Für einen score-basierten Klassifikator (z.\,B.\ $h_\theta(x)$ oder die SVM-
Marginweite) mit Schwellwert liefert die \textbf{Konfusionsmatrix} vier
Grundgrößen: \textbf{TP} (true positive), \textbf{TN}, \textbf{FP}
("Fehler 1. Art"), \textbf{FN} ("Fehler 2. Art"). Daraus:

\begin{center}
\begin{tabular}{lll}
\toprule
Metrik & Formel & Interpretation \\
\midrule
Accuracy & $\dfrac{TP+TN}{TP+TN+FP+FN}$ & Anteil korrekter Vorhersagen (= 1 -- 0/1-Verlust) \\
Precision & $\dfrac{TP}{TP+FP}$ & Anteil echter Positiver unter den vorhergesagten Positiven \\
Recall (Sensitivity) & $\dfrac{TP}{TP+FN}$ & Anteil gefundener echter Positiver \\
Specificity & $\dfrac{TN}{TN+FP}$ & Anteil gefundener echter Negativer \\
F1-Score & $\dfrac{2\cdot\text{Prec}\cdot\text{Rec}}{\text{Prec}+\text{Rec}}$ & Harmonisches Mittel aus Precision und Recall \\
\bottomrule
\end{tabular}
\end{center}

\begin{tcolorbox}[defstyle]
\textbf{Precision-Recall-Tradeoff.} Ein niedrigerer Schwellwert erhöht
Recall (mehr werden als positiv vorhergesagt), senkt aber typischerweise
Precision -- und umgekehrt. \emph{Triviale} $100\,\%$ Recall: alle Beispiele
als positiv vorhersagen. \emph{Triviale} $100\,\%$ Precision: nur das
sicherste Beispiel als positiv vorhersagen. Der F1-Score fasst diesen
Tradeoff in einer Zahl zusammen.
\end{tcolorbox}

\subsection{Summary-Metriken: ROC- und PR-Kurven}

Statt einen einzelnen Schwellwert zu fixieren, variiert man ihn über den
gesamten Wertebereich und erhält Kurven:

\begin{itemize}
  \item \textbf{ROC-Kurve} (Receiver Operating Characteristic): Sensitivity
        gegen $(1-\text{Specificity})$. \textbf{AUROC} = Fläche darunter =
        $P(\text{zufälliges Positiv-Beispiel ranked höher als zufälliges
        Negativ-Beispiel})$. Unabhängig von der Klassenprävalenz.
  \item \textbf{PR-Kurve} (Precision-Recall): Precision gegen Recall.
        \textbf{AUPRC} = erwartete Precision bei zufälligem Schwellwert;
        bei starkem Klassenungleichgewicht aussagekräftiger als AUROC.
\end{itemize}

\begin{tcolorbox}[defstyle]
\textbf{Klassenungleichgewicht (Prävalenz $<5\,\%$).} Accuracy wird trivial
hoch, indem stets die Mehrheitsklasse vorhergesagt wird (Baseline =
Prävalenz). AUROC bleibt leicht hoch haltbar, indem die meisten Negativen
sehr niedrig gescort werden. AUPRC ist robuster gegen dieses Verhalten,
aber nicht immun. Faustregel: $\text{Accuracy} \prec \text{AUROC} \prec
\text{AUPRC}$ in Robustheit gegenüber starkem Klassenungleichgewicht.
\end{tcolorbox}

\textbf{Log-Loss} (siehe Kreuzentropie, Abschnitt~2.3) und der eng verwandte
\textbf{Brier-Score} $\frac1n\sum_i(\hat p_i - y_i)^2$ bestrafen -- anders
als Accuracy/AUROC/AUPRC, die nur auf dem \emph{Ranking} beruhen -- explizit
schlecht kalibrierte, übermäßig selbstbewusste Fehlvorhersagen und sind
\emph{proper scoring rules}: minimal genau dann, wenn die vorhergesagten
Wahrscheinlichkeiten den wahren entsprechen.

% ==================================================================
\newpage
\part{Unüberwachtes Lernen}
% ==================================================================

% ------------------------------------------------------------------
\section{K-Means-Clustering}
% ------------------------------------------------------------------

K-Means partitioniert $n$ Datenpunkte in $K$ Cluster durch Minimierung der
quadratischen Intra-Cluster-Distanz:
\begin{equation}
  J(c, \mu) = \sum_{i=1}^n \norm{x^{(i)} - \mu_{c^{(i)}}}^2.
\end{equation}

\begin{tcolorbox}[algostyle, title={K-Means Algorithmus}]
\begin{enumerate}[leftmargin=*, itemsep=2pt]
  \item Initialisiere $K$ Zentroide $\mu_1,\ldots,\mu_K$ (z.\,B.\ zufällig).
  \item Wiederhole bis Konvergenz:
    \begin{enumerate}[label=(\alph*)]
      \item \textbf{Assignment}: $c^{(i)} \leftarrow \arg\min_k \norm{x^{(i)}-\mu_k}^2$
      \item \textbf{Update}: $\mu_k \leftarrow \frac{\sum_i \I\{c^{(i)}=k\}\,x^{(i)}}
              {\sum_i \I\{c^{(i)}=k\}}$
    \end{enumerate}
\end{enumerate}
\end{tcolorbox}

Jede Iteration verringert $J$; der Algorithmus konvergiert zu einem
(möglicherweise lokalen) Minimum. \textbf{K-Means++}-Initialisierung
(proportional zu $\norm{x - \mu_\text{nächster}}^2$) verbessert die
Lösungsqualität systematisch.

% ------------------------------------------------------------------
\section{EM-Algorithmus}
% ------------------------------------------------------------------

\subsection{Mixture of Gaussians (MoG)}

Ein Gaußsches Gemischmodell nimmt an, dass jeder Datenpunkt $x^{(i)}$ von
einem der $K$ Gaußschmischkomponenten erzeugt wurde:
\begin{equation}
  p(x^{(i)}) = \sum_{k=1}^K \phi_k\,\N(x^{(i)};\,\mu_k, \Sigma_k),
\end{equation}
wobei $\phi_k \geq 0$, $\sum_k \phi_k = 1$. Durch latente Variablen $z^{(i)}
\in \{1,\ldots,K\}$ (Clusterzugehörigkeit) wird das Modell:
$z^{(i)} \sim \text{Multinomial}(\phi)$, $x^{(i)} \mid z^{(i)}=k \sim
\N(\mu_k, \Sigma_k)$.

\subsection{Jensens Ungleichung}

\begin{tcolorbox}[thmstyle]
\textbf{Jensens Ungleichung.}
Für eine konkave Funktion $f$ und eine Zufallsvariable $X$:
$\E[f(X)] \leq f(\E[X])$.
Für konvexes $f$: $\E[f(X)] \geq f(\E[X])$.
\end{tcolorbox}

Dies ist die Basis für die ELBO-Schranke: mit $Q$ als beliebiger Verteilung
über $z$:
\begin{equation}
  \log p(x;\,\theta)
  \geq \sum_z Q(z)\log\frac{p(x,z;\,\theta)}{Q(z)}
  =: \mathrm{ELBO}(Q, \theta).
\end{equation}

\subsection{Allgemeiner EM-Algorithmus}

\begin{tcolorbox}[algostyle, title={Expectation-Maximization (EM)}]
Wiederhole bis Konvergenz:
\begin{itemize}[leftmargin=*, itemsep=2pt]
  \item \textbf{E-Schritt}: $Q_i(z^{(i)}) \leftarrow p(z^{(i)} \mid x^{(i)};\,\theta)$
  \item \textbf{M-Schritt}:
    $\theta \leftarrow \arg\max_\theta \sum_i \sum_{z^{(i)}} Q_i(z^{(i)}) \log\frac{p(x^{(i)}, z^{(i)};\,\theta)}{Q_i(z^{(i)})}$
\end{itemize}
\end{tcolorbox}

\textbf{Konvergenz}: Jede Iteration erhöht $\log p(x;\,\theta)$ oder lässt
es gleich. Die ELBO-Schranke ist beim E-Schritt scharf (Gleichheit in
Jensens Ungleichung).

\subsection{EM für Gauß-Gemische}

\textbf{E-Schritt}: Berechne Zugehörigkeitswahrscheinlichkeiten
(\emph{responsibilities}):
\begin{equation}
  \gamma_{ik} = Q_i(z^{(i)}=k)
  = \frac{\phi_k\,\N(x^{(i)};\,\mu_k,\Sigma_k)}
         {\sum_{j=1}^K \phi_j\,\N(x^{(i)};\,\mu_j,\Sigma_j)}.
\end{equation}

\textbf{M-Schritt}: Update der Parameter:
\begin{align}
  \hat\phi_k &= \frac{1}{n}\sum_i \gamma_{ik}, \\
  \hat\mu_k  &= \frac{\sum_i \gamma_{ik}\, x^{(i)}}{\sum_i \gamma_{ik}}, \\
  \hat\Sigma_k &= \frac{\sum_i \gamma_{ik}(x^{(i)}-\hat\mu_k)(x^{(i)}-\hat\mu_k)\T}
                       {\sum_i \gamma_{ik}}.
\end{align}

% ------------------------------------------------------------------
\section{Faktorenanalyse}
% ------------------------------------------------------------------

\subsection{Motivation: Wenn $n \gg m$}

Bei hochdimensionalen Daten ($n \gg m$, mehr Features als Trainingsbeispiele)
ist die MLE-Kovarianzschätzung
$\hat\Sigma = \tfrac1m\sum_i (x^{(i)}-\hat\mu)(x^{(i)}-\hat\mu)\T$ singulär
(Rang $\leq m-1 < n$), sodass $\hat\Sigma^{-1}$ nicht existiert -- die
Gauß-Dichte ist damit für GDA oder Gauß-Gemische nicht auswertbar.
Einschränkungen wie diagonales $\Sigma$ oder $\Sigma=\sigma^2 I$ vermeiden
Singularität, modellieren dann aber keinerlei Korrelationen zwischen
Features mehr. Die \textbf{Faktorenanalyse} bietet einen Kompromiss: ein
niedrigdimensionales latentes Modell, das dennoch Korrelationsstruktur
einfängt, ohne eine volle $n\times n$-Kovarianzmatrix zu schätzen.

\subsection{Das Modell}

Mit latenter Variable $z \in \R^k$ ($k \ll n$):
\begin{equation}
  z \sim \N(0, I), \qquad x \mid z \sim \N(\mu + \Lambda z,\, \Psi),
\end{equation}
äquivalent $x = \mu + \Lambda z + \varepsilon$ mit $\varepsilon\sim\N(0,\Psi)$
und $z \indep \varepsilon$. Die Parameter sind $\mu\in\R^n$,
$\Lambda\in\R^{n\times k}$ (Faktor-Ladungsmatrix) und die \emph{diagonale}
Matrix $\Psi\in\R^{n\times n}$ (features-spezifisches Rauschen).

Da $(z,x)$ gemeinsam Gauß-verteilt sind, ergibt sich die Randverteilung
\begin{equation}
  x \sim \N(\mu,\; \Lambda\Lambda\T + \Psi),
\end{equation}
d.\,h.\ die Kovarianz von $x$ wird durch $\Lambda\Lambda\T$ (Rang $\leq k$,
erfasst die Korrelationen über gemeinsame latente Faktoren) plus diagonalem
Rauschen $\Psi$ approximiert -- mit nur $O(nk)$ statt $O(n^2)$ freien
Parametern.

\subsection{EM für Faktorenanalyse}

Da die Log-Likelihood in $(\mu,\Lambda,\Psi)$ nicht geschlossen maximierbar
ist, wird EM eingesetzt (vgl.\ Abschnitt zum EM-Algorithmus).

\textbf{E-Schritt}: Über die Formeln für bedingte Gauß-Verteilungen
($\mu_{1|2} = \mu_1+\Sigma_{12}\Sigma_{22}^{-1}(x_2-\mu_2)$,
$\Sigma_{1|2}=\Sigma_{11}-\Sigma_{12}\Sigma_{22}^{-1}\Sigma_{21}$, angewendet
auf die Gemeinsamverteilung von $z$ und $x$) ergibt sich
\begin{align}
  z^{(i)} \mid x^{(i)};\,\mu,\Lambda,\Psi &\sim \N\bigl(\mu_{z^{(i)}|x^{(i)}},\, \Sigma_{z^{(i)}|x^{(i)}}\bigr), \\
  \mu_{z^{(i)}|x^{(i)}} &= \Lambda\T(\Lambda\Lambda\T+\Psi)^{-1}(x^{(i)}-\mu), \\
  \Sigma_{z^{(i)}|x^{(i)}} &= I - \Lambda\T(\Lambda\Lambda\T+\Psi)^{-1}\Lambda.
\end{align}

\textbf{M-Schritt}: Update von $\Lambda$ und $\Psi$ unter Nutzung der
Momente $\E[z^{(i)}\mid x^{(i)}]$ und $\E[z^{(i)}(z^{(i)})\T\mid x^{(i)}]$ aus
dem E-Schritt (geschlossene, aber längere Update-Formeln analog zum
Gauß-Gemisch-Fall).

\begin{tcolorbox}[defstyle]
\textbf{Bezug zu PCA und ICA.} Faktorenanalyse, PCA (nächster Abschnitt) und
ICA suchen alle eine niedrigdimensionale latente Repräsentation von $x$.
Faktorenanalyse ist ein \emph{probabilistisches} generatives Modell
(explizites Rauschen $\Psi$, EM-Schätzung); PCA entspricht dem Grenzfall
$\Psi = \sigma^2 I$ für $\sigma^2\to 0$ und liefert dafür eine \emph{direkte},
nicht-iterative Lösung über eine Eigenwertzerlegung.
\end{tcolorbox}

% ------------------------------------------------------------------
\section{Hidden Markov Models}
% ------------------------------------------------------------------

\begin{tcolorbox}[srcstyle]
Quelle: \emph{Hidden Markov Models Fundamentals}, Daniel Ramage, CS229
Section Notes -- \texttt{section/cs229-hmm.pdf}.
\end{tcolorbox}

Faktorenanalyse modelliert eine latente Variable \emph{pro Beispiel};
Hidden Markov Models modellieren eine latente Variable \emph{pro Zeitschritt}
in einer Sequenz -- und werden, wie Faktorenanalyse und Gauß-Gemische, per
EM (hier "Baum-Welch" genannt) gelernt.

\subsection{Markov-Modelle}

Gegeben Zustände $S=\{s_1,\ldots,s_{|S|}\}$ und eine beobachtete
Zustandsfolge $z_1,\ldots,z_T \in S$. Zwei Annahmen machen das Problem
handhabbar:
\begin{itemize}
  \item \textbf{Limited-Horizon-Annahme}: $P(z_t \mid z_{t-1},\ldots,z_1) = P(z_t\mid z_{t-1})$.
  \item \textbf{Stationaritätsannahme}: $P(z_t\mid z_{t-1})$ ist unabhängig von $t$.
\end{itemize}
Mit Übergangsmatrix $A_{ij} = P(z_t{=}s_j \mid z_{t-1}{=}s_i)$ und einem
künstlichen Startzustand $z_0$ ergibt sich die Sequenzwahrscheinlichkeit per
Kettenregel als
\begin{equation}
  P(\vec z; A) = \prod_{t=1}^T P(z_t \mid z_{t-1}; A) = \prod_{t=1}^T A_{z_{t-1} z_t}.
\end{equation}

\subsection{Hidden Markov Models: Drei Fragen}

Bei einem \textbf{HMM} sind die Zustände $z_1,\ldots,z_T$ \emph{nicht}
beobachtbar -- man beobachtet nur eine Ausgabe $x_t\in V$ pro Zeitschritt,
erzeugt gemäß der \textbf{Output-Independence-Annahme}
$P(x_t{=}v_k\mid z_t{=}s_j) = B_{jk}$ (Emissionsmatrix $B$). Drei
Grundfragen:

\begin{enumerate}
  \item \textbf{Likelihood}: Wie wahrscheinlich ist eine Beobachtungsfolge $\vec x$?
  \item \textbf{Decoding}: Welche Zustandsfolge $\vec z$ erklärt $\vec x$ am besten?
  \item \textbf{Learning}: Wie schätzt man $A,B$ aus Beobachtungen?
\end{enumerate}

\subsection{Forward-Prozedur (Frage 1)}

Naiv erfordert $P(\vec x;A,B) = \sum_{\vec z} P(\vec x\mid\vec z)P(\vec z)$
eine Summe über $|S|^T$ Zustandsfolgen. Mit der Hilfsgröße
$\alpha_i(t) = P(x_1,\ldots,x_t,\,z_t{=}s_i;A,B)$ liefert dynamische
Programmierung eine $O(|S|^2 T)$-Lösung:

\begin{tcolorbox}[algostyle, title={Forward-Prozedur}]
\begin{enumerate}[leftmargin=*, itemsep=1pt]
  \item Basisfall: $\alpha_i(0) = A_{0i}$, $i=1,\ldots,|S|$.
  \item Rekursion: $\alpha_j(t) = \sum_{i=1}^{|S|} \alpha_i(t{-}1)\,A_{ij}\,B_{j x_t}$, $j=1,\ldots,|S|$, $t=1,\ldots,T$.
  \item $P(\vec x;A,B) = \sum_{i=1}^{|S|} \alpha_i(T)$.
\end{enumerate}
\end{tcolorbox}

Analog liefert die \textbf{Backward-Prozedur}
$\beta_i(t) = P(x_{t+1},\ldots,x_T \mid z_t{=}s_i;A,B)$.

\subsection{Viterbi-Algorithmus (Frage 2)}

Gesucht: $\arg\max_{\vec z} P(\vec z\mid \vec x;A,B) = \arg\max_{\vec z}
P(\vec x,\vec z;A,B)$. Der \textbf{Viterbi-Algorithmus} ist strukturell
identisch zur Forward-Prozedur, ersetzt aber die Summe $\sum_i$ durch ein
Maximum $\max_i$ und speichert bei jedem Schritt zusätzlich den
Vorgängerzustand -- so lässt sich am Ende die optimale Zustandsfolge per
Backtracking rekonstruieren. Laufzeit ebenfalls $O(|S|^2 T)$.

\subsection{Baum-Welch-Algorithmus (Frage 3)}

Parameterlernen für HMMs ist ein Spezialfall des EM-Algorithmus
(Abschnitt~16). Mit der Hilfsgröße
$\gamma_t(i,j) = \alpha_i(t)\,A_{ij}\,B_{jx_{t+1}}\,\beta_j(t{+}1)$
(proportional zur Wahrscheinlichkeit eines Übergangs $s_i\to s_j$ zur Zeit
$t$ gegeben alle Beobachtungen):

\begin{tcolorbox}[algostyle, title={Baum-Welch (Forward-Backward-EM)}]
Wiederhole bis Konvergenz:
\begin{itemize}[leftmargin=*, itemsep=1pt]
  \item \textbf{E-Schritt}: Forward/Backward laufen lassen, $\gamma_t(i,j)$ berechnen.
  \item \textbf{M-Schritt}:
    \[
      A_{ij} \leftarrow \frac{\sum_{t=1}^T \gamma_t(i,j)}{\sum_{j=1}^{|S|}\sum_{t=1}^T \gamma_t(i,j)}, \qquad
      B_{jk} \leftarrow \frac{\sum_{i=1}^{|S|}\sum_{t:\,x_t=v_k} \gamma_t(i,j)}{\sum_{i=1}^{|S|}\sum_{t=1}^T \gamma_t(i,j)}.
    \]
\end{itemize}
\end{tcolorbox}

Intuitiv: $A_{ij}$ = erwartete Anzahl Übergänge $s_i\to s_j$, geteilt durch
erwartete Häufigkeit von $s_i$; $B_{jk}$ analog für Emissionen. Wie jede
EM-Anwendung ist das Problem nicht-konvex (mehrere lokale Maxima) --
mehrere Initialisierungen und Smoothing (keine Nullwahrscheinlichkeiten)
sind ratsam.

% ------------------------------------------------------------------
\section{Hauptkomponentenanalyse (PCA)}
% ------------------------------------------------------------------

PCA findet die Richtungen größter Varianz in den Daten und projiziert auf
den $k$-dimensionalen Unterraum, der die erklärte Varianz maximiert.

\textbf{Algorithmus:}
\begin{enumerate}
  \item Zentriere Daten: $x^{(i)} \leftarrow x^{(i)} - \bar x$.
  \item Berechne Stichproben-Kovarianzmatrix: $\hat\Sigma = \frac{1}{n}X\T X$.
  \item Eigenwertzerlegung: $\hat\Sigma = U \Lambda U\T$ (oder SVD von $X$).
  \item Wähle die $k$ Eigenvektoren $u_1,\ldots,u_k$ zu den größten
        Eigenwerten.
  \item Projektion: $z^{(i)} = U_k\T x^{(i)} \in \R^k$.
\end{enumerate}

\textbf{Interpretation:} $u_j$ ist die $j$-te Hauptkomponente.
Der Anteil der erklärten Varianz durch $k$ Komponenten ist
$\sum_{j=1}^k \lambda_j / \sum_{j=1}^d \lambda_j$.

PCA ist die Lösung des Optimierungsproblems
\begin{equation}
  \max_{U\T U = I_k} \operatorname{tr}(U\T \hat\Sigma U).
\end{equation}

% ------------------------------------------------------------------
\section{Unabhängige Komponentenanalyse (ICA)}
% ------------------------------------------------------------------

ICA löst das \textbf{Blind Source Separation}-Problem: Beobachtungen
$x = As$ sind lineare Mischungen unabhängiger Quellen $s \in \R^d$.
Ziel ist die Schätzung der Mischmatrix $A$ ohne Kenntnis von $s$.

\subsection{ICA-Ambiguitäten}

Grundlegende Nichtidentifizierbarkeiten:
\begin{itemize}
  \item \textbf{Skalierung}: $As = (A\,\mathrm{diag}(c))\,(\mathrm{diag}(1/c)\,s)$
        -- Quellen können nur bis auf einen Skalierungsfaktor bestimmt werden.
  \item \textbf{Permutation}: Reihenfolge der Quellen ist unbestimmt.
  \item \textbf{Gaussianität}: Gaußverteilte Quellen sind \emph{nie}
        identifizierbar (alle Richtungen äquivalent).
\end{itemize}

\subsection{ICA-Algorithmus}

Sei $W = A^{-1}$ die Entmischungsmatrix. Die Likelihood der Beobachtungen ist:
\begin{equation}
  p_x(x) = \prod_{j=1}^d p_s(w_j\T x) \cdot |\det W|.
\end{equation}
Der ICA-Algorithmus maximiert die Log-Likelihood durch stochastischen
Gradient Ascent auf $W$:
\begin{equation}
  W \leftarrow W + \alpha
  \Bigl(\bigl[1 - 2\,g(Wx)\bigr]\,(Wx)\T + (W\T)^{-1}\Bigr),
\end{equation}
wobei $g(s) = 1/(1+e^{-s})$ die Sigmoid-Funktion ist (unter der Annahme
einer Super-Gaußschen Quellenverteilung).

% ------------------------------------------------------------------
\section{Selbstüberwachtes Lernen und Foundation Models}
% ------------------------------------------------------------------

\subsection{Motivation}

Manuelle Annotation großer Datensätze ist teuer. \textbf{Selbstüberwachtes
Lernen} nutzt in den Daten selbst enthaltene Strukturen als Supervision
-- ohne menschliche Labels.

\subsection{Pretraining in Computer Vision}

\begin{itemize}
  \item \textbf{Kontrastives Lernen} (SimCLR, MoCo): Lernt Repräsentationen,
        indem zwei Augmentierungen desselben Bildes als ``positiv'' und
        verschiedene Bilder als ``negativ'' behandelt werden.
  \item \textbf{Masked Image Modeling} (MAE): Lernt, maskierte Bildregionen
        aus dem Kontext zu rekonstruieren.
\end{itemize}

\subsection{Large Language Models und Transformer}

\textbf{Transformer-Architektur:}
\begin{itemize}
  \item \textbf{Self-Attention}: $\mathrm{Attn}(Q,K,V) = \mathrm{softmax}(QK\T/\sqrt{d_k})\,V$.
        Lernt, welche Positionen im Kontext relevant sind.
  \item \textbf{Multi-Head Attention}: Mehrere parallele Attention-Heads
        erfassen verschiedene linguistische Aspekte.
  \item \textbf{Feedforward-Schicht}: Positionsweise angewendetes MLP nach
        jeder Attention-Schicht.
  \item \textbf{Positionskodierung}: Addiert Information über die Position
        des Tokens, da Attention permutationsinvariant ist.
\end{itemize}

\textbf{Pretraining}: Das Modell lernt, das nächste Token zu
vorherzusagen:
\[
  \mathcal{L} = -\sum_t \log P(x_t \mid x_1,\ldots,x_{t-1};\,\theta).
\]

\textbf{Zero-Shot- und In-Context-Learning}: Vortrainierte LLMs können neue
Aufgaben ohne Fine-Tuning lösen, wenn die Aufgabenbeschreibung (und ggf.\
Beispiele) im Prompt mitgegeben werden -- ein emergentes Phänomen großer
Modelle.

% ------------------------------------------------------------------
\section{Kritische Perspektiven auf ML}
% ------------------------------------------------------------------

\begin{tcolorbox}[srcstyle]
Quelle: \emph{Technical and Societal Critiques of ML}, Chris Chute, Taide
Ding, Andrey Kurenkov, CS229 Fall 2019 --
\texttt{materials/critiques-ml-aut19.pdf} (vgl.\ auch die Vorgängerversion
\texttt{materials/critiques-ml.pdf}).
\end{tcolorbox}

Leistungsfähige Modelle sind nicht automatisch vertrauenswürdige oder
robuste Modelle. Drei Schwachstellen, die über reine Testfehler-Metriken
hinausgehen:

\subsection{Adversarial Examples}

Neuronale Netze verletzen oft die intuitive Annahme, dass sich die
Modellvorhersage in einer kleinen Umgebung $\norm{r}<\varepsilon$ um ein
Trainingsbeispiel $x$ kaum ändert. Der \textbf{Fast Gradient Sign Method
(FGSM)}-Angriff konstruiert eine minimale, für Menschen kaum wahrnehmbare
Störung, die eine falsche Klassifikation erzwingt, indem er in Richtung des
\emph{steigenden} Verlusts geht:
\begin{equation}
  x_\text{adv} = x + \varepsilon\cdot\operatorname{sign}\bigl(\nabla_x J(\theta,x,y)\bigr).
\end{equation}
Solche Störungen generalisieren häufig über Modellarchitekturen und
Trainingsdatensätze hinweg und lassen sich teils sogar physisch realisieren
(z.\,B.\ präparierte Verkehrsschilder). Verteidigungen (adversariales
Training, Distillation) mindern das Problem, gelten aber weiterhin als
offene Forschungsfrage.

\subsection{Interpretierbarkeit}

Nach Lipton ("The Mythos of Model Interpretability") lassen sich zwei
Perspektiven unterscheiden:
\begin{itemize}
  \item \textbf{Algorithmische Transparenz}: Das Modell ist selbst
        zerlegbar/verständlich (Beispiel: Koeffizienten der linearen
        Regression). Klassische lineare Modelle und einzelne Decision Trees
        gewinnen hier.
  \item \textbf{Post-hoc-Erklärung}: Nachträgliche Erklärungen via
        Visualisierung, Saliency Maps, Grad-CAM oder Attention-Gewichte.
        Hier können auch tiefe neuronale Netze reichhaltige Erklärungen
        liefern.
\end{itemize}
Die verbreitete Vereinfachung "lineare Modelle = interpretierbar, neuronale
Netze = Blackbox" verwechselt diese beiden Perspektiven. Kernfrage hinter
Interpretierbarkeit: Vertrauen, Kausalität, Übertragbarkeit auf neue
Verteilungen und Fairness -- Aspekte, die reine Test-Genauigkeits- oder
AUROC-Metriken (Abschnitt~16) nicht erfassen.

\subsection{Kosten und gesellschaftliche Aspekte}

Ergänzend diskutiert werden in den Folien der hohe Daten- und
Rechenaufwand moderner Modelle sowie ethische Fragestellungen (Bias in
Trainingsdaten, Fairness der Modellentscheidungen gegenüber
unterrepräsentierten Gruppen) -- Themen, die bei der praktischen
Anwendung von ML (nicht zuletzt bei Foundation Models, Abschnitt~23)
zunehmend an Bedeutung gewinnen.

% ==================================================================
\newpage
\part{Bestärkendes Lernen}
% ==================================================================

% ------------------------------------------------------------------
\section{Bestärkendes Lernen}
% ------------------------------------------------------------------

\subsection{Markov-Entscheidungsprozesse}

Ein \textbf{Markov-Entscheidungsprozess (MDP)} ist ein Tupel
$(S, A, P, R, \gamma)$:
\begin{itemize}
  \item $S$: Zustandsraum, $A$: Aktionenraum,
  \item $P_{sa}(s') = P(s' \mid s, a)$: Übergangsdynamik,
  \item $R : S \times A \to \R$: Belohnungsfunktion,
  \item $\gamma \in [0,1)$: Diskontierungsfaktor.
\end{itemize}

Ziel: Lerne eine \textbf{Policy} $\pi : S \to A$, die den erwarteten
diskontierten Gesamtertrag maximiert:
\begin{equation}
  V^\pi(s) = \E\!\Bigl[\sum_{t=0}^\infty \gamma^t R(s_t, a_t)
    \;\Big|\; s_0 = s,\, \pi\Bigr].
\end{equation}

\subsection{Bellman-Gleichungen}

Die \textbf{optimale Wertfunktion} $V^*$ erfüllt die
\textbf{Bellman-Optimalitätsgleichung}:
\begin{equation}
  V^*(s) = \max_a \Bigl[R(s,a) + \gamma \sum_{s'} P_{sa}(s')\,V^*(s')\Bigr].
\end{equation}

Die optimale Policy ist $\pi^*(s) = \arg\max_a [R(s,a) + \gamma \sum_{s'}
P_{sa}(s') V^*(s')]$.

\subsection{Value Iteration und Policy Iteration}

\begin{tcolorbox}[algostyle, title={Value Iteration}]
Initialisiere $V(s) = 0$. Wiederhole bis Konvergenz:
\[
  V(s) \leftarrow \max_a \Bigl[R(s,a) + \gamma \sum_{s'} P_{sa}(s')\,V(s')\Bigr]
  \quad \forall s
\]
\end{tcolorbox}

\begin{tcolorbox}[algostyle, title={Policy Iteration}]
Initialisiere Policy $\pi$ beliebig. Wiederhole:
\begin{enumerate}[leftmargin=*, itemsep=1pt]
  \item \textbf{Policy Evaluation}: Löse linea\-res Gleichungssystem
        $(I - \gamma P^\pi)V^\pi = R^\pi$ nach $V^\pi$.
  \item \textbf{Policy Improvement}:
        $\pi(s) \leftarrow \arg\max_a [R(s,a) + \gamma \sum_{s'} P_{sa}(s') V^\pi(s')]$
\end{enumerate}
Policy Iteration konvergiert in endlich vielen Schritten.
\end{tcolorbox}

\subsection{Wertfunktions-Approximation (kontinuierliche Zustände)}

Bei kontinuierlichem Zustandsraum $S \subseteq \R^d$ approximiert man
$V(s) \approx \theta\T \phi(s)$ und aktualisiert per Fitted Value Iteration:
\begin{equation}
  y^{(i)} = R(s^{(i)}) + \gamma \max_a \sum_{s'} P_{s^{(i)} a}(s')\,V(s'),
  \quad \theta \leftarrow \arg\min_\theta \sum_i (\theta\T\phi(s^{(i)}) - y^{(i)})^2.
\end{equation}

% ------------------------------------------------------------------
\section{Lineare Quadratische Regelung, DDP und LQG}
% ------------------------------------------------------------------

\subsection{Lineare Quadratische Regelung (LQR)}

LQR behandelt Systeme mit \textbf{linearer Dynamik} und
\textbf{quadratischen Kosten}:
\begin{align}
  x_{t+1} &= Ax_t + Bu_t + w_t, \quad w_t \sim \N(0, \Sigma_w), \\
  \text{Kosten:} &\quad c_t(x_t, u_t) = x_t\T Q\, x_t + u_t\T R\, u_t,
  \quad Q \succeq 0,\, R \succ 0.
\end{align}

Die optimale Wertfunktion ist quadratisch: $V_t(x) = x\T P_t x + q_t$.
Die \textbf{Riccati-Rekursion} liefert $P_t$ rückwärts in der Zeit:
\begin{align}
  P_T &= Q_T, \\
  P_t &= Q + A\T P_{t+1} A
    - A\T P_{t+1} B(R + B\T P_{t+1} B)^{-1} B\T P_{t+1} A.
\end{align}
Die optimale Policy ist \textbf{linear im Zustand}:
\begin{equation}
  u_t^* = -K_t x_t, \qquad
  K_t = (R + B\T P_{t+1} B)^{-1} B\T P_{t+1} A.
\end{equation}

\subsection{Differential Dynamic Programming (DDP)}

DDP erweitert LQR auf \textbf{nichtlineare Systeme}
$x_{t+1} = f(x_t, u_t)$ durch lokale Linearisierung um eine
Referenztrajektorie $(\bar x_t, \bar u_t)$:
\begin{equation}
  f(x,u) \approx f(\bar x,\bar u) + F_x(x - \bar x) + F_u(u - \bar u).
\end{equation}
Backward Pass (wie LQR) berechnet die Feedback-Gains; Forward Pass wendet
die Policy an und aktualisiert die Referenztrajektorie. DDP iteriert bis
Konvergenz.

\subsection{Lineares Quadratisches Gauß-Modell (LQG)}

LQG kombiniert LQR mit einem \textbf{Kalman-Filter}, wenn der Zustand $x_t$
nicht direkt beobachtbar ist:
\begin{equation}
  y_t = Cx_t + v_t, \quad v_t \sim \N(0, \Sigma_v).
\end{equation}

\textbf{Separationsprinzip:} Kalman-Filter und LQR-Regler können unabhängig
voneinander entworfen werden -- der Kalman-Filter schätzt $\hat x_t = \E[x_t
\mid y_1,\ldots,y_t]$, der LQR-Regler nutzt $\hat x_t$ als wäre es der wahre
Zustand.

% ------------------------------------------------------------------
\section{Policy Gradient (REINFORCE)}
% ------------------------------------------------------------------

\subsection{Direkte Policy-Parametrisierung}

Statt eine Wertfunktion zu lernen, parametrisiert man die Policy direkt:
$\pi_\theta(a \mid s)$. Das Optimierungsziel ist:
\begin{equation}
  J(\theta) = \E_\tau\!\Bigl[\sum_{t=0}^T \gamma^t r_t\Bigr],
  \quad \tau \sim \pi_\theta.
\end{equation}

\subsection{Policy-Gradient-Theorem}

\begin{tcolorbox}[thmstyle]
\textbf{Policy-Gradient-Theorem.}
\[
  \nabla_\theta J(\theta)
  = \E_\tau\!\Bigl[\sum_{t=0}^T \nabla_\theta \log\pi_\theta(a_t \mid s_t)
    \cdot G_t\Bigr],
\]
wobei $G_t = \sum_{t'=t}^T \gamma^{t'-t} r_{t'}$ der zukünftige diskontierte
Ertrag ist.
\end{tcolorbox}

Das Ergebnis nutzt den Trick
$\nabla_\theta \pi_\theta = \pi_\theta \nabla_\theta \log\pi_\theta$ und
die REINFORCE-Identität.

\subsection{REINFORCE-Algorithmus}

\begin{tcolorbox}[algostyle, title={REINFORCE}]
\begin{enumerate}[leftmargin=*, itemsep=2pt]
  \item Sampling: Führe Episode $\tau = (s_0, a_0, r_0, \ldots, s_T)$ unter
        $\pi_\theta$ aus.
  \item Berechne $G_t = \sum_{t'=t}^T \gamma^{t'-t} r_{t'}$.
  \item Gradient: $g = \sum_t \nabla_\theta \log\pi_\theta(a_t \mid s_t) \cdot G_t$.
  \item $\theta \leftarrow \theta + \alpha\, g$.
\end{enumerate}
\end{tcolorbox}

\subsection{Varianzreduktion}

REINFORCE hat hohe Varianz. \textbf{Baselines} reduzieren sie ohne den
Erwartungswert zu verzerren:
\begin{equation}
  g = \sum_t \nabla_\theta \log\pi_\theta(a_t \mid s_t)
    \cdot \bigl(G_t - b(s_t)\bigr).
\end{equation}
Typische Wahl: $b(s_t) = V^\pi(s_t)$. Die Differenz
$A(s,a) = Q(s,a) - V(s)$ heißt \textbf{Vorteilsfunktion (Advantage)} und
führt auf \textbf{Actor-Critic}-Methoden.

% ==================================================================
\newpage
\section*{Schluss: Zusammenfassung der Schlüsselkonzepte}
% ==================================================================

\addcontentsline{toc}{section}{Schluss: Zusammenfassung der Schlüsselkonzepte}

Die folgende Tabelle gibt einen Überblick über die wichtigsten Algorithmen
und ihre Eigenschaften:

\begin{center}
\begin{tabular}{p{3.2cm} p{3cm} p{3.2cm} p{4.5cm}}
\toprule
\textbf{Methode} & \textbf{Aufgabe} & \textbf{Schlüsselidee} & \textbf{Stärken / Grenzen} \\
\midrule
Lineare Regression & Regression & Least Squares & Einfach, gut interpretierbar; linear \\
Logist. Regression & Binäre Klassif. & Sigmoid + MLE & Probabilistisch; benötigt Linearität \\
GDA & Klassifikation & Gauß-Modell $p(x|y)$ & Gut bei wenig Daten; Gauß-Annahme \\
Naive Bayes & Klassifikation & Bed. Unabhängigkeit & Skalierbar, einfach; starke Annahme \\
SVM & Klassifikation & Maximaler Margin & Kernel-fähig; teuer bei großem $n$ \\
Perceptron & Klassifikation & Online-Update bei Fehler & Mistake-Bound $(D/\gamma)^2$; nur bei Trennbarkeit \\
Decision Tree & Beides & Rekursive Regionssplits & Interpretierbar; hohe Varianz \\
Random Forest & Beides & Bagging + Feature-Subsampling & Robust, wenig Tuning; weniger interpretierbar \\
Boosting (AdaBoost) & Klassifikation & Sequentielle schwache Lerner & Reduziert Bias; overfitting-anfällig bei großem $M$ \\
Neuronale Netze & Allgemein & Tiefe Nichtlinearität & Universell; viele Hyperparameter \\
K-Means & Clustering & Centroid-Update & Einfach; lokale Optima, wählt $K$ \\
EM & Dichtemodell & E/M-Schritt & Allgemein; langsam, lokale Optima \\
Faktorenanalyse & Dim.-Reduktion & Latentes Gauß-Modell + EM & Gut bei $n \gg m$; iterativ \\
Hidden Markov Model & Sequenzen & Forward/Viterbi/Baum-Welch & Zeitreihen mit latentem Zustand; nicht-konvex \\
Gaussian Process & Regression & Bayes-Prior über Funktionen & Kalibrierte Unsicherheit; teuer bei großem $n$ ($O(n^3)$) \\
PCA & Dim.-Reduktion & Eigenvektoren & Linear; verliert nichtlineare Struktur \\
Value Iteration & RL (diskret) & Bellman-Update & Exakt; skaliert schlecht \\
LQR & RL (kontinuierl.) & Riccati, linear & Analytisch; erfordert lineare Dynamik \\
REINFORCE & RL (allgemein) & Policy Gradient & Flexibel; hohe Varianz \\
\bottomrule
\end{tabular}
\end{center}

\bigskip

\noindent Diese Vorlesung führt von einfachen parametrischen Modellen über
tiefe neuronale Netze bis hin zu bestärkendem Lernen für sequentielle
Entscheidungen. Das verbindende Thema ist stets die Optimierung einer
Verlust- oder Wertfunktion -- meist durch Gradientenverfahren, manchmal
durch geschlossene Formeln oder dynamische Programmierung.

\bigskip

\noindent\textit{``Machine learning is the science of getting computers
to act without being explicitly programmed.''} --- Andrew Ng

\end{document}
